% Pass table option to xcolor BEFORE aaai2027 loads it (for \rowcolor support)
\PassOptionsToPackage{table}{xcolor}
\documentclass[letterpaper]{article} % DO NOT CHANGE THIS
\usepackage[submission]{aaai2027}  % DO NOT CHANGE THIS
% Fonts (newtxtext, helvet, courier) are loaded automatically by aaai2027.sty.
% DO NOT add \usepackage{times}, \usepackage{helvet}, or \usepackage{courier}.
\usepackage[hyphens]{url}  % DO NOT CHANGE THIS
\usepackage{graphicx} % DO NOT CHANGE THIS
\urlstyle{rm} % DO NOT CHANGE THIS
\def\UrlFont{\rm}  % DO NOT CHANGE THIS
\usepackage{natbib}  % DO NOT CHANGE THIS AND DO NOT ADD ANY OPTIONS TO IT
\usepackage{caption} % DO NOT CHANGE THIS AND DO NOT ADD ANY OPTIONS TO IT
\frenchspacing  % DO NOT CHANGE THIS
\setlength{\pdfpagewidth}{8.5in} % DO NOT CHANGE THIS
\setlength{\pdfpageheight}{11in} % DO NOT CHANGE THIS

\pdfinfo{
/TemplateVersion (2027.1)
}

\usepackage{booktabs}
\usepackage{amsfonts}
\usepackage{nicefrac}
\usepackage{microtype}
\usepackage{amsmath}
\usepackage{amssymb}
\usepackage{subcaption}
\usepackage{multirow}
\usepackage{siunitx}
\usepackage{enumitem}
\usepackage{makecell}
\usepackage{threeparttable}
\usepackage{longtable}
\title{EgoEMG: A Multimodal Egocentric Dataset with Bilateral EMG and Vision for Hand Pose Estimation}

\author{
    Anonymous Authors
}
\affiliations{
    Anonymous Institution\\
    anonymous@institution.edu
}

\newcommand{\dataset}{\textit{EgoEMG}}
\newcommand{\model}{EMGFormer}
\newcommand{\marker}{\textit{markers2mano}}
\newcommand{\degree}{^\circ}
% Define \mainonly at the command line to build the 9-page AAAI main submission.
\newif\ifmainonly
\ifdefined\mainonly
  \mainonlytrue
\else
  \mainonlyfalse
\fi
\newcommand{\supplementref}[1]{%
  the supplementary material}
\newcommand{\supplementfigureref}[1]{%
  \ifmainonly supplementary material\else Figure~\ref{#1}\fi}
\begin{document}

\maketitle

\begin{abstract}
Surface electromyography (sEMG) captures muscle activity during hand movement, while egocentric vision observes global hand configuration. These complementary signals are promising for robust hand-pose estimation, but no existing dataset jointly provides synchronized bilateral wrist EMG, egocentric vision, and bimanual pose labels. We present EgoEMG, a multimodal egocentric dataset containing bilateral wristband EMG sampled at 2\,kHz, 120\,Hz IMU signals, egocentric video, external RGB-D video, and motion-capture-derived wrist and finger joint angles. EgoEMG comprises over 18 hours of recordings from 53 participants performing 60 gesture classes---30 unimanual and 30 bimanual. We establish benchmarks for EMG-to-pose, vision-to-pose, and EMG+vision fusion under shared joint-angle targets and Gesture, User, and Both generalization splits. Across seven visual backbones, residual fusion reduces average error relative to the corresponding vision-only baseline. Together, EgoEMG provides a foundation for studying multimodal hand-pose estimation from muscle activity and visual observations.
\end{abstract}

\section{Introduction}

Hand pose estimation from wearable sensing is a building block for augmented reality, prosthetic control, and physical intelligence~\citep{emg2pose,ninapro,zhao2026dexemg}. Surface EMG is attractive because it measures muscle activation directly~\citep{farina2014} and remains usable under self-occlusion, motion blur, and challenging lighting. Recent work has shown that EMG can recover detailed hand articulation~\citep{emg2pose,emg2tendon,vqmyopose}. At the same time, egocentric vision has emerged as a powerful modality for hand tracking~\citep{hamer,h2o}, but relies on line-of-sight and struggles with fast motion blur and self-occlusion during manipulation. These properties suggest a potential complementarity between the two modalities: EMG provides fine-grained finger articulation even when the hand is occluded, while vision recovers global hand configuration and spatial context. Unlocking this complementarity requires a dataset that synchronizes both modalities with pose labels---yet current EMG pose datasets provide no synchronized egocentric video.

We address these gaps with \dataset{}, a controlled multimodal egocentric dataset for bimanual hand pose estimation. \dataset{} contains synchronized bilateral wristband EMG (16 total channels, 8 per wrist at 2\,kHz), per-hand IMU, egocentric wide-angle RGB, external RGB-D, and hand motion reconstructed into MANO and UmeTrack parameters~\citep{mano,umetrack} with wrist articulation angles, covering 53 participants and 60 gesture classes (30 single-hand, 30 bimanual) over 18\,hours of recording.

We provide a benchmark with three tasks---EMG-to-pose, vision-to-pose, and EMG+vision fusion---sharing the same joint-angle prediction target and Gesture, User, and Both generalization splits. On the EMG2Pose dataset, \model{}-L reduces User+Stage MAE by 16\% relative to the best prior reported result. On \dataset{}, \model{}-M achieves the lowest overall MAE among all evaluated EMG baselines. On vision-to-pose, ResNet backbones outperform ViT at lower parameter counts, while the hand-specialized WiLoR baseline achieves the strongest standalone visual accuracy. Across all seven visual backbones, residual fusion lowers average error relative to the matched vision-only baseline, and the gains become more modest for the strongest visual models.

Our contributions are:
\begin{itemize}[leftmargin=1.5em,nosep]
    \item \dataset{}, the first dataset jointly providing bilateral wristband EMG and IMU, egocentric RGB, external RGB-D, bimanual MANO/UmeTrack annotations with wrist joint angles, and per-frame gesture class labels.
    \item A three-task hand pose estimation benchmark (EMG-to-pose, vision-to-pose, EMG+vision fusion) with Gesture, User, and Both generalization splits under a unified joint-angle prediction target.
    \item Baseline results showing that \model{} provides a strong EMG-to-pose reference on both EMG2Pose and \dataset{}, that hand-specialized vision modeling improves over generic ResNet/ViT baselines, and that residual EMG+vision fusion improves accuracy across the evaluated visual backbones.
\end{itemize}  

\section{Related Work}

\paragraph{Hand pose datasets and benchmarks.}
Public EMG datasets cover gesture recognition~\citep{ninapro,grabmyo,putemg,du2017surface,amma2015advancing}, high-density recordings~\citep{jiang2021open}, force estimation~\citep{pimforce}, typing~\citep{emg2qwerty}, and demographically diverse cohorts~\citep{gowda2025database}. Recent work also combines wrist-worn sEMG and IMU with egocentric video to obtain force-enriched manipulation demonstrations~\citep{forceband}. For continuous pose estimation, EMG2Pose~\citep{emg2pose} provides synchronized bilateral EMG with mocap-derived joint angles, while MyoKi~\citep{myoki} and \citet{jarque2019} pair forearm EMG with instrumented-glove kinematics. Cross-user generalization has been studied through benchmarks such as EMGBench~\citep{yang2024emgbench}. In contrast, visual hand-pose benchmarks provide rich RGB annotations for monocular, egocentric, and hand-object settings~\citep{freihand,dexycb,interhand,hot3d,h2o,gigahand}, but do not include synchronized wrist EMG. Existing resources therefore lack a unified benchmark that combines egocentric video, bilateral wristband EMG, wrist articulation, and continuous hand-pose labels, limiting the study of multimodal pose regression from both visual observations and muscle activity.

\paragraph{EMG-based pose estimation.}
Decoding hand kinematics from surface EMG has progressed from single-DoF~\citep{ngeo2012continuous,ngeo2014continuous} and multi-DoF finger-angle regression to full-hand pose estimation. NeuroPose~\citep{liu2021neuropose} demonstrated 3D hand-pose tracking from wearable EMG, while EMG2Pose~\citep{emg2pose} later established a larger-scale benchmark using a TDS+LSTM architecture. Recent work explores tendon-driven motion modeling~\citep{emg2tendon}, vector-quantized pose tokenization~\citep{vqmyopose}, short-calibration user adaptation~\citep{react2026}, the role of position versus velocity decoding under the original benchmark protocol~\citep{hadidi2026reevaluating}, and generalization-focused teleoperation~\citep{zhao2026dexemg}. In parallel, EMG foundation models and self-supervised approaches~\citep{fasulo2025tinymyo,mehlman2025scaling,vicregemg,mast} explore large-scale pretraining, but primarily target gesture classification rather than continuous pose. Despite these advances, EMG-only pose estimation remains sensitive to inter-user variability, and datasets based on instrumented gloves~\citep{myoki,jarque2019} provide limited degrees of freedom and can suffer from sensor drift and hysteresis.

\paragraph{Vision-based hand pose estimation and multimodal fusion.}
Monocular RGB hand-pose estimation has advanced rapidly with the MANO hand model~\citep{mano}, large-scale annotated datasets~\citep{freihand,dexycb,interhand,gigahand}, and Transformer-based reconstructors~\citep{hamer,wilor}. Egocentric systems and datasets such as UmeTrack~\citep{umetrack}, HOT3D~\citep{hot3d}, H2O~\citep{h2o}, and WristPP~\citep{xi2026wristpp} further support first-person hand tracking and interaction understanding in wearable and egocentric settings. However, vision-based methods remain fundamentally constrained by line of sight: self-occlusion, inter-hand occlusion, object manipulation, and motion blur can substantially degrade pose estimates. EMG provides complementary muscle-activation cues that are independent of image visibility. Recent work has begun exploring EMG+vision multimodal learning for gesture understanding~\citep{hao2024multimodal,zandigohar2024multimodal,embridge} and cross-modal EMG-pose representation~\citep{cpep,liu2025frompose}, but these efforts focus on gesture classification or offline representation mapping rather than continuous hand-pose regression. The lack of synchronized egocentric video, bilateral EMG, and continuous pose labels has therefore prevented systematic study of EMG+vision fusion for hand-pose regression.

\section{\dataset{} Dataset}
\label{sec:dataset}

\subsection{Overview and Capture Setup}

\dataset{} is collected with a synchronized multimodal capture setup, including
bilateral EMG, IMU, egocentric RGB, external RGB-D, and optical motion capture.
Each participant wears two WAVELETECH EMG wristbands, one on each wrist, with each
wristband recording 8 channels of surface EMG at 2\,kHz. The visual setup
contains a head-mounted wide-angle RGB camera for egocentric hand observation
and an external ZED 2i RGB-D camera mounted on the motion-capture rack. Hand
motion is tracked using 21 reflective markers per hand at 120\,Hz with an FZMotion optical motion capture system. All modalities
are temporally synchronized via host-timestamp soft-synchronization
followed by linear interpolation to a common timeline (Figure~\ref{fig:capture_setup}).

We recruited 53 participants. \dataset{} contains 60 gesture classes (30 single-hand, 30 bimanual) with per-frame gesture labels for gesture-conditioned evaluation. Gesture prompts are displayed in randomized order; single-hand gestures are performed mirror-symmetrically with both hands while varying wrist orientation; bimanual gestures follow demonstrated two-hand interactions, capturing natural coordination patterns. Figure~\ref{fig:sample_data} shows representative synchronized samples from all modalities.

\begin{figure}[t]
\centering
\includegraphics[width=0.9\columnwidth]{figures/aaai/fig_capture_setup.pdf}
\caption{Data collection setup for \dataset{}. Bilateral EMG wristbands, head-mounted egocentric RGB, external ZED 2i RGB-D, and optical motion capture with hand markers for pose labels.}
\label{fig:capture_setup}
\end{figure}

\begin{figure}[t]
\centering
\includegraphics[width=0.95\columnwidth]{figures/aaai/fig_sample_data.pdf}
\caption{Representative synchronized samples from \dataset{}. Each row (6\,s window) shows GT hand pose, bilateral EMG, egocentric RGB at window center, and external RGB-D.}
\label{fig:sample_data}
\end{figure}

\subsection{Pose Label Reconstruction}
\label{subsec:label_reconstruction}

Hand pose labels are reconstructed from optical motion-capture markers into
MANO parameters~\citep{mano}. We use a learning-based \marker{} pipeline that
maps 21 marker positions to MANO pose and shape parameters. The
pipeline uses fixed MANO surface vertices as virtual marker correspondences,
allowing captured 3D marker positions to be aligned with the canonical MANO
mesh. Full details of the \marker{} architecture, training data, and augmentation
pipeline are provided in \supplementref{app:markers2mano}.

Compared with the per-frame inverse-kinematics solver used by EMG2Pose, which
has a reported 12.7\% invalid-frame rate, our learning-based reconstruction
reduces the invalid-frame rate to 3.6\%, a 3.5$\times$ reduction. The resulting
MANO meshes achieve a mean marker-to-mesh alignment error of 4.3\,mm.
In addition to MANO parameters, \dataset{} provides per-frame joint angles in a
22-DoF scalar format (20 finger angles plus 2 wrist articulation angles) obtained
by optimizing UmeTrack joint angles to match MANO landmark positions;
the conversion procedure is described in
\supplementref{app:joint_angle_conversion}.
As an independent cross-modal check, we reproject reconstructed MANO meshes
onto egocentric RGB frames using the calibrated camera model. Visual inspection of randomly sampled
reprojections confirms consistent mesh-to-image alignment across gesture
families (\supplementfigureref{fig:label_quality}).

Wrist articulation angles are computed directly from motion-capture markers,
independent of the MANO reconstruction pipeline. Markers attached to each EMG
armband define a forearm coordinate frame, from which we compute two wrist
degrees of freedom: flexion/extension and radial/ulnar deviation; the full
derivation is provided in \supplementref{app:wrist_angles}.

\subsection{Dataset Statistics and Positioning}

Table~\ref{tab:dataset_compare} compares \dataset{} with representative public
EMG datasets. Among existing datasets, only EMG2Pose and \dataset{} provide
bilateral EMG with motion-capture-derived continuous hand-pose labels.
\dataset{} further contributes synchronized egocentric RGB, external RGB-D,
IMU, and explicit wrist articulation labels, enabling direct study of
cross-modal hand pose estimation.

\begin{table}[t]
\centering
\caption{Comparison of \dataset{} with representative public EMG datasets.
``EMG Ch.'' counts total channels across both wrists (per-wrist$\times$2 for bilateral datasets); ``Bilateral'' indicates dual-wrist EMG capture; ``Wrist labels'' indicates explicit wrist articulation labels; ``Pose labels'' indicates continuous hand pose annotations.}
\label{tab:dataset_compare}
\resizebox{0.9\columnwidth}{!}{%
\begin{threeparttable}
\begin{tabular}{@{}lcccccccc@{}}
\toprule
Dataset & Subj. & EMG Ch. & Bilateral & Ego RGB & RGB-D &
\makecell{Wrist\\labels} & \makecell{Pose\\labels} & Gestures \\
\midrule
Jarque-Bou et al.~\shortcite{jarque2019} & 22 & 8 & -- & -- & -- & -- & 18 (g) & 26 \\
NinaPro~\citep{ninapro} & 27--78 & 10--16 & -- & -- & -- & \checkmark & 22 (g)\tnote{*} & 53 \\
MyoKi~\citep{myoki} & 36 & 16 & -- & -- & -- & -- & 18 (g) & 74 \\
PiMForce~\citep{pimforce} & 21 & 8 & -- & -- & -- & -- & 20 (g) & 63 \\
EMG2Pose~\citep{emg2pose} & 193 & 32 (16$\times$2) & \checkmark & -- & -- & -- & 20 (m) & 50 \\
\textbf{\dataset{} (ours)}
& \textbf{53} & \textbf{16 (8$\times$2)} & \checkmark & \checkmark & \checkmark
& \checkmark & \textbf{22 (m)} & \textbf{60} \\
\bottomrule
\end{tabular}
\begin{tablenotes}
\item[*] NinaPro DB1, 2 and 5 provide continuous joint angles (22-DoF); remaining sub-databases provide gesture labels only.
\item[] (g) = instrumented glove; (m) = optical motion capture.
\end{tablenotes}
\end{threeparttable}
}
\end{table}

\paragraph{Hand shape and gesture diversity.}
Figure~\ref{fig:dataset_analysis}(a) shows the MANO shape distribution
across 53 participants, demonstrating substantial inter-subject hand-shape
variation. The 60 gesture classes include 30 single-hand gestures and 30
bimanual interaction gestures, covering both symmetric coordination and
asymmetric manipulation. The full gesture vocabulary is provided in
\supplementref{app:gesture_vocabulary}.

\begin{figure}[t]
    \centering
    \includegraphics[width=\columnwidth]{figures/aaai/statistics.pdf}
    \caption{\textbf{(a)} MANO shape diversity across \dataset{} participants (106 hands). \textbf{(b)} Evaluation splits on the participant--gesture matrix (Train, Gesture/User/Both).}
    \label{fig:dataset_analysis}
\end{figure}

\subsection{Evaluation Splits and Release Format}

Following the protocol of EMG2Pose~\citep{emg2pose}, we define two primary
generalization axes: \emph{gesture} and \emph{user}. Based on these axes, we
evaluate three mutually exclusive splits: a \emph{Gesture split}, where held-out gestures are
evaluated on seen participants; a \emph{User split}, where held-out participants
are evaluated on seen gestures; and a combined \emph{Both split}, where both
gestures and participants are held out. All modalities follow the same split
assignment, preserving synchronized multimodal examples and preventing leakage
along the held-out axis. We hold out 10 gestures for the gesture split and
6 participants for the user split, with an overall train-to-validation/test
ratio of approximately 7:3. Figure~\ref{fig:dataset_analysis}(b) illustrates the
split design.

The released dataset includes synchronized EMG and IMU signals, egocentric RGB
frames, external RGB-D frames, MANO-based and UmeTrack-based hand-pose
annotations, wrist articulation angles, timestamps, and calibration metadata.
The prediction target is a 22-DoF joint-angle vector per hand: 20 finger angles
and 2 wrist articulation angles. The same target definition and semantic ordering
are used for both left and right hands.

\section{Benchmark and Baselines}

\subsection{Benchmark Tasks}

As an initial benchmark, \dataset{} defines three prediction tasks under the same 22-DoF target
definition and evaluation split axes. The target consists of 20 finger joint angles
and 2 wrist articulation angles per hand. The three tasks are:
\emph{EMG-to-pose}, which predicts joint angles from EMG windows;
\emph{vision-to-pose}, which predicts joint angles from single egocentric RGB hand
crops; and \emph{EMG+vision fusion}, which uses both modalities to predict the
same output. The external RGB-D and IMU streams are included in the release to
support future extensions.

\subsection{Evaluation Protocol}

We evaluate all tasks on the Gesture, User, and Both splits defined in the
dataset section above. The primary metric is mean absolute error (MAE) over joint
angles:
\begin{equation}
\text{MAE} =
\frac{1}{JT}
\sum_{j=1}^{J}\sum_{t=1}^{T}
|\hat{\theta}_{j,t} - \theta_{j,t}|
\quad [^\circ],
\end{equation}
where $J$ is the number of predicted joint angles and $T$ is the number of
evaluated time steps. For wrist analysis, we
additionally report MAE on flexion/extension and radial/ulnar deviation. Unless
otherwise specified, all errors are reported on the held-out test split. For
vision-only and fusion evaluation, we use $256\times256$ hand crops localized by
projecting motion-capture keypoints into the calibrated egocentric camera. These
models predict the window center frame; fusion pairs that image with a centered,
non-causal 6\,s EMG window. The resulting protocol evaluates offline pose
regression given hand localization rather than hand detection or causal tracking.
EMG-only evaluation instead uses the full predicted trajectory over each window.

\subsection{EMG Baseline}

As the EMG-to-pose baseline, we introduce \model{} (Figure~\ref{fig:architecture}),
a temporal model that maps filtered EMG windows to per-frame joint angles. For
the reported \dataset{} experiments, each example contains 12{,}000 samples
(6\,s) from the target wrist's eight EMG channels. A TDS-style temporal
convolutional frontend~\citep{tds} downsamples the 2\,kHz signal to latent
features at approximately 38\,Hz, followed by a Transformer
decoder with rotary positional encoding (RoPE)~\citep{rope} to model long-range
temporal dependencies. Left-hand EMG is mirrored to the right-hand
convention during training. Each forward pass predicts one hand's joint-angle
trajectory, including wrist articulation for \dataset{}. We provide three model
sizes; implementation details, architecture specifications, and
ablations are provided in \supplementref{app:training_details}.

\begin{figure}[t]
    \centering
    \includegraphics[width=0.9\columnwidth]{figures/aaai/emg.pdf}
    \caption{\model{} architecture for EMG-to-pose. Target-wrist EMG is encoded
    by a TDS-style temporal convolutional frontend, decoded by a Transformer with
    RoPE positional encoding, and projected to joint angles including wrist
    articulation.}
    \label{fig:architecture}
\end{figure}

\subsection{Vision Baseline}

For vision-to-pose, we evaluate seven single-frame reference baselines. Six are generic image backbones across two families: ResNet~\citep{resnet} (\{18, 50, 152\}) and ViT~\citep{vit} (\{Small, Base, Large\}). ResNet backbones are initialized from ImageNet~\citep{imagenet}-pretrained weights, while ViT backbones are initialized from DINOv2~\citep{dinov2}-pretrained weights. Each generic model takes a
$256\times256$ egocentric hand crop from the window center frame, extracts a
feature vector with the backbone, and regresses the 22-DoF target through an MLP
prediction head. We additionally evaluate WiLoR~\citep{wilor} as a hand-specialized vision baseline initialized from a pretrained MANO-based hand-pose checkpoint and fine-tuned on the same protocol (Figure~\ref{fig:vision_architecture}). Model
sizes and training configurations are provided in \supplementref{app:training_details_vision}.

\begin{figure}[t]
    \centering
    \includegraphics[width=0.9\columnwidth]{figures/aaai/vision.pdf}
    \caption{Vision-to-pose baseline. A $256\times256$ hand crop from the
    center frame is processed by a ResNet or ViT backbone, followed by an MLP
    head that regresses 22 joint angles.}
    \label{fig:vision_architecture}
\end{figure}

\subsection{Fusion Baseline}
\label{sec:fusion_arch}

The EMG+vision baseline is designed as a simple residual fusion model
(Figure~\ref{fig:fusion_architecture}). The vision branch first predicts a
vision-only pose $\mathbf{y}_v$, while the EMG branch predicts an additive
correction $\Delta\mathbf{y}_{\text{emg}}$. The final prediction is $\mathbf{y} = \mathbf{y}_v + \Delta\mathbf{y}_{\text{emg}}$.
This design makes the fusion model start from a strong visual baseline and
learn the incremental pose information contributed by EMG.

The vision branch extracts a 256-d feature from the center-frame hand crop and
uses a prediction head to produce $\mathbf{y}_v \in \mathbb{R}^{22}$. The EMG
branch uses \model{} to encode the full EMG window into a temporal feature
sequence, followed by learned temporal attention pooling to obtain a 256-d EMG
feature. The two features are concatenated and passed to a fusion head that
predicts $\Delta\mathbf{y}_{\text{emg}} \in \mathbb{R}^{22}$. The last layer of
the residual head is zero-initialized so that the model begins from the
vision-only prediction. We train the fusion model under a center-supervised
protocol: only the window center frame is supervised, matching the vision-only
target and enabling direct comparison.

\begin{figure}[t]
    \centering
    \includegraphics[width=0.9\columnwidth]{figures/aaai/fusion.pdf}
    \caption{EMG+vision fusion baseline. Vision branch predicts $\mathbf{y}_v$; EMG branch predicts residual $\Delta\mathbf{y}_{\text{emg}}$ (zero-initialized head).}
    \label{fig:fusion_architecture}
\end{figure}

\section{Experiments}

\subsection{EMG Baseline Results}

Table~\ref{tab:egoemg_emg_results} reports EMG-to-pose results on \dataset{} for \model{} variants and four traditional architectures across the Gesture, User, and Both test conditions. The baselines are the original emg2pose TDS+MLP model and the recurrent vemg2pose model~\citep{emg2pose}, NeuroPose, a convolutional encoder-decoder~\citep{liu2021neuropose}, and SensingDynamics, a convolutional model operating on raw and low-pass EMG~\citep{sensingdynamics}.

\begin{table}[t]
\centering
\caption{EMG-to-pose results on \dataset{}. All baselines are retrained under
the same evaluation protocol. MAE is reported in degrees; lower is better.
Gesture/User/Both columns report per-user mean $\pm$ standard deviation. Avg
is the per-sample weighted mean across all test splits. The shaded row marks
the best accuracy--capacity trade-off; bold marks the lowest unambiguous result
in each column.}
\label{tab:egoemg_emg_results}
\footnotesize
\setlength{\tabcolsep}{3.8pt}
\renewcommand{\arraystretch}{1.08}
\resizebox{\columnwidth}{!}{%
\begin{tabular}{lccccc}
\toprule
Method & Params & Gesture & User & Both & Avg. \\
\midrule
\multicolumn{6}{l}{\textit{\model{} variants}} \\
\midrule
\model{}-S & 3.5M & $12.3_{\pm1.5}$ & $16.0_{\pm0.6}$ & $16.3_{\pm1.5}$ & $14.1$ \\
\rowcolor{gray!12}
\model{}-M & 6.6M & $\mathbf{11.7_{\pm1.6}}$ & $15.9_{\pm0.6}$ & $16.4_{\pm1.5}$ & $\mathbf{13.8}$ \\
\model{}-L & 16.3M & $11.9_{\pm1.6}$ & $16.0_{\pm0.8}$ & $16.4_{\pm1.2}$ & $13.9$ \\
\midrule
\multicolumn{6}{l}{\textit{Reimplemented baselines}} \\
\midrule
emg2pose~\citep{emg2pose} & 3.0M & $15.5_{\pm1.3}$ & $\mathbf{14.8_{\pm2.9}}$ & $16.3_{\pm0.6}$ & $15.8$ \\
vemg2pose~\citep{emg2pose} & 6.0M & $15.0_{\pm1.4}$ & $16.3_{\pm1.7}$ & $17.3_{\pm1.3}$ & $15.9$ \\
NeuroPose~\citep{liu2021neuropose} & 6.4M & $15.8_{\pm1.2}$ & $15.7_{\pm1.3}$ & $16.3_{\pm0.7}$ & $16.1$ \\
SensingDynamics~\citep{sensingdynamics} & 1.0M & $16.2_{\pm1.1}$ & $16.4_{\pm0.3}$ & $16.7_{\pm0.8}$ & $16.4$ \\
\bottomrule
\end{tabular}}%

\end{table}

\model{} variants achieve the best accuracy on the Gesture split, indicating that \model{} better captures gesture-level articulation under the seen-user, held-out-gesture setting. The original emg2pose baseline retains the lowest User-split error, underscoring the challenge of cross-user generalization, while \model{}-S matches the lowest reported Both-split error. \model{}-M achieves the best overall MAE and offers the best accuracy--capacity trade-off among the \model{} variants, outperforming \model{}-L with fewer than half the parameters.

\begin{table}[t]
\centering
\caption{EMG-only regression results on EMG2Pose test splits. MAE is reported in degrees; lower is better. \model{} values report per-user mean $\pm$ standard deviation across users. $^\dagger$ indicates numbers reported in the corresponding papers. MT denotes joint Regression+Tracking training, evaluated here in the Regression setting without test-time user calibration. Bold marks the lowest result in each column.}
\label{tab:emg2pose_results}
\footnotesize
\setlength{\tabcolsep}{3.5pt}
\renewcommand{\arraystretch}{1.05}
\resizebox{\columnwidth}{!}{%
\begin{tabular}{lcccc}
\toprule
Method & Params & User & Stage & User+Stage \\
\midrule
\multicolumn{5}{l}{\textit{\model{} variants}} \\
\midrule
\model{}-S & 3.5M & $12.5_{\pm1.1}$ & $11.1_{\pm1.2}$ & $12.3_{\pm1.1}$ \\
\model{}-M & 6.6M & $12.4_{\pm1.1}$ & $10.2_{\pm1.1}$ & $12.4_{\pm1.1}$ \\
\model{}-L & 16.3M & $12.3_{\pm1.1}$ & $\mathbf{9.3_{\pm1.1}}$ & $\mathbf{12.3_{\pm1.1}}$ \\
\midrule
\multicolumn{5}{l}{\textit{Prior work (reported)}} \\
\midrule
emg2pose (position)$^\dagger$~\citep{emg2pose} & 3.0M & $12.6_{\pm1.3}$ & $15.2_{\pm1.6}$ & $15.6_{\pm1.3}$ \\
vemg2pose$^\dagger$~\citep{emg2pose} & 6.0M & $12.2_{\pm1.3}$ & $15.2_{\pm1.6}$ & $15.8_{\pm1.4}$ \\
NeuroPose$^\dagger$~\citep{liu2021neuropose} & 6.4M & $13.2_{\pm1.1}$ & $17.2_{\pm1.7}$ & $17.5_{\pm1.5}$ \\
SensingDynamics$^\dagger$~\citep{sensingdynamics} & 1.0M & $15.5_{\pm1.4}$ & $18.8_{\pm1.6}$ & $18.7_{\pm1.6}$ \\
Position MT$^\dagger$~\citep{hadidi2026reevaluating} & 6.0M & $11.5_{\pm1.2}$ & $14.0_{\pm1.6}$ & $14.6_{\pm1.3}$ \\
Velocity MT$^\dagger$~\citep{hadidi2026reevaluating} & 6.0M & $11.6_{\pm1.3}$ & $13.9_{\pm1.6}$ & $14.6_{\pm1.3}$ \\
CLDM$^\dagger$~\citep{emg2tendon} & 7.0M & $\mathbf{11.3_{\pm1.0}}$ & $14.3_{\pm1.5}$ & $14.7_{\pm1.4}$ \\
\bottomrule
\end{tabular}}%
\end{table}

On EMG2Pose (Table~\ref{tab:emg2pose_results}), \model{}-L achieves 12.25$^\circ$ $\pm$ 1.08$^\circ$ MAE on User+Stage, a 16\% reduction relative to the best prior reported result of 14.6$^\circ$ $\pm$ 1.3$^\circ$ from Position/Velocity MT~\citep{hadidi2026reevaluating}. The compact \model{}-S reaches 12.34$^\circ$ $\pm$ 1.07$^\circ$ with only 58\% of the vemg2pose model size.

\subsection{Vision-to-Pose Results}

Table~\ref{tab:egoemg_results} reports vision-only and EMG+vision results on
\dataset{}. We first analyze the vision-only rows, which include six generic
single-frame baselines spanning two backbone families, ResNet and ViT, together
with one hand-specialized baseline, WiLoR. Among the generic backbones,
ResNet-152 achieves the best vision-only performance across all splits.
ResNet-50 also outperforms ViT-Large despite using substantially fewer
parameters, suggesting that convolutional backbones remain highly competitive as
generic egocentric references in this data regime. Scaling reduces error within
both generic backbone families, while the User and Both splits remain harder
than the Gesture split across all vision-only models. The relative
underperformance of generic ViT regressors may reflect the scale of the
available supervised training data, as transformer-based vision models often
benefit from larger datasets. In contrast, the hand-specialized WiLoR baseline
achieves the strongest standalone visual accuracy, reducing the average MAE to
4.7$^\circ$. This indicates that task-specific hand-pose pretraining remains important in the \dataset{} setting.

\subsection{EMG+Vision Fusion Results}
\label{sec:cross_modal}

We evaluate EMG+vision fusion using the residual architecture described above,
with \model{}-S as the EMG backbone. The fusion model predicts the window
center frame, where the egocentric RGB observation is available, making its MAE
directly comparable to the vision-only baseline under the same supervision
target. For each visual backbone, we train two fusion variants: one freezes the
pretrained visual predictor and the other fine-tunes it jointly with the EMG
branch. For each backbone, we retain the frozen or fine-tuned variant with the
lower validation MAE. Table~\ref{tab:egoemg_results} reports the retained
variant; the complete frozen-versus-fine-tuned comparison is deferred to
\supplementref{app:fusion_freeze_ablation}.

\begin{table*}[t]
\centering
\caption{Vision-only versus EMG+vision fusion on
\dataset{}. Each split reports MAE in degrees as Vision $\rightarrow$ Fusion on
identical center frames; bold marks the lower error. Frz./FT denotes a
frozen/fine-tuned visual predictor. Params reports total parameters in millions
as vision-only/fusion. Avg is the mean and
$\Delta_{\mathrm{avg}}$ is the overall fusion gain computed from unrounded
values.}
\label{tab:egoemg_results}
\footnotesize
\setlength{\tabcolsep}{4.8pt}
\renewcommand{\arraystretch}{1.08}
\begin{tabular}{lccccccc}
\toprule
Backbone & Params (V/F) & Update & Gesture & User & Both & Avg & $\Delta_{\mathrm{avg}}$ \\
\midrule
ResNet-18  & 11.5/15.5M  & FT & $5.10\!\rightarrow\!\mathbf{4.60}$ & $6.80\!\rightarrow\!\mathbf{6.40}$ & $6.80\!\rightarrow\!\mathbf{6.40}$ & $5.84\!\rightarrow\!\mathbf{5.40}$ & $+0.44$ \\
ResNet-50  & 23.5/29.0M  & Frz. & $4.50\!\rightarrow\!\mathbf{4.39}$ & $6.23\!\rightarrow\!\mathbf{6.19}$ & $6.10\!\rightarrow\!\mathbf{6.05}$ & $5.27\!\rightarrow\!\mathbf{5.19}$ & $+0.09$ \\
ResNet-152 & 58.2/63.6M  & Frz. & $4.31\!\rightarrow\!\mathbf{4.25}$ & $6.07\!\rightarrow\!\mathbf{6.04}$ & $6.03\!\rightarrow\!6.03$ & $5.11\!\rightarrow\!\mathbf{5.06}$ & $+0.05$ \\
\midrule
ViT-S/14   & 21.9/25.8M  & FT & $5.00\!\rightarrow\!\mathbf{4.50}$ & $7.30\!\rightarrow\!\mathbf{6.80}$ & $7.20\!\rightarrow\!\mathbf{6.80}$ & $6.02\!\rightarrow\!\mathbf{5.54}$ & $+0.48$ \\
ViT-B/14   & 86.2/90.2M  & Frz. & $4.85\!\rightarrow\!\mathbf{4.80}$ & $\mathbf{7.02}\!\rightarrow\!7.03$ & $6.86\!\rightarrow\!\mathbf{6.84}$ & $5.78\!\rightarrow\!\mathbf{5.75}$ & $+0.03$ \\
ViT-L/14   & 303.8/308.0M & Frz. & $4.39\!\rightarrow\!\mathbf{4.34}$ & $\mathbf{6.67}\!\rightarrow\!6.69$ & $6.60\!\rightarrow\!6.60$ & $5.39\!\rightarrow\!\mathbf{5.36}$ & $+0.03$ \\
\midrule
WiLoR      & 631.6/635.8M & Frz. & $3.98\!\rightarrow\!\mathbf{3.96}$ & $5.57\!\rightarrow\!\mathbf{5.52}$ & $5.58\!\rightarrow\!\mathbf{5.54}$ & $4.73\!\rightarrow\!\mathbf{4.68}$ & $+0.04$ \\
\bottomrule
\end{tabular}

\end{table*}

Fusion lowers the sample-weighted average error for all seven backbones, with the largest gains for ResNet-18 and ViT-S and smaller gains for stronger visual models. The smaller gains with larger visual backbones may reflect the growing parameter imbalance between the visual backbone and the much smaller EMG and fusion branches, which may make joint multimodal optimization more difficult. The improvements are not uniform across splits: frozen ViT-B and ViT-L improve on unseen gestures but are nearly unchanged or slightly worse on unseen users, whereas WiLoR yields small gains on all three splits. The ablation in \supplementref{app:fusion_freeze_ablation} indicates that fine-tuning works better for ResNet-18 and ViT-S, while freezing is preferable for larger backbones under our settings. The per-gesture analysis in \supplementref{app:per_gesture} further examines where fusion is most beneficial.

\paragraph{Effect of self-occlusion.}
We analyze the matched Vision-RN18 and Fusion-RN18+S models. To quantify self-occlusion, we rasterize each ground-truth MANO mesh in camera coordinates at the native image resolution using the calibrated camera and a z-buffer. A vertex is visible when its depth agrees with the local z-buffer within 5\,mm. With $w_i$ denoting the surface-area weight of vertex $i$, we define
\begin{equation}
s_{\mathrm{occ}} = 1 - \frac{\sum_{i \in \mathcal{V}_{\mathrm{vis}}} w_i}
{\sum_{i=1}^{V} w_i},
\end{equation}
where $\mathcal{V}_{\mathrm{vis}}$ is the visible-vertex set; $s_{\mathrm{occ}}=0$ denotes a fully visible mesh (full details in \supplementref{app:occlusion}). Fusion improves on vision-only in every occlusion bin: the mean gain rises from $0.24\degree$ in the least-occluded bin to $0.52\degree$ in the most-occluded bin, and reaches approximately $0.8\degree$ at moderate-to-high occlusion. Figure~\ref{fig:occlusion_gain} visualizes this pattern. Although the frame-level association between occlusion and fusion gain is weak, the positive gains across bins are consistent with EMG providing complementary cues when visual evidence is degraded. Figure~\ref{fig:qualitative_fusion} illustrates a representative case: when visual evidence is ambiguous, EMG retains pose cues independent of image appearance, and the fused model recovers articulation closer to the ground truth.

\begin{figure}[t]
  \centering
  \includegraphics[width=\columnwidth]{figures/occlusion_fusion_gain.png}
  \caption{Fusion benefit as a function of hand self-occlusion on the \dataset{} test set (16{,}650 valid predictions). \textbf{(a)} Mean joint-angle MAE of vision-only and EMG+vision fusion across fixed-width (0.05) occlusion bins; annotations show the fusion gain at the lowest and highest bins. \textbf{(b)} Per-bin fusion gain ($\Delta = \text{MAE}_{\text{vision}} - \text{MAE}_{\text{fusion}}$) with a LOESS trend line. Error bands denote $\pm$1 SEM.}
  \label{fig:occlusion_gain}
\end{figure}


\begin{figure}[t]
\centering
\includegraphics[width=0.95\columnwidth]{figures/quality.png}
\caption{Qualitative comparison on \dataset{}. The fused model (green) recovers articulation closer to GT under motion blur, self-occlusion, and depth ambiguity.}
\label{fig:qualitative_fusion}
\end{figure}

\section{Discussion and Conclusion}

\dataset{} provides a multimodal benchmark for hand pose estimation with
synchronized bilateral EMG, IMU, egocentric RGB, external RGB-D, and hand-pose
annotations including wrist articulation. By using a shared 22-DoF prediction
target and unified split axes, the benchmark aligns EMG-to-pose,
vision-to-pose, and EMG+vision fusion under a common label space while
retaining modality-specific supervision protocols.

Our experiments highlight three findings. First, cross-user generalization
remains the central challenge for EMG-based pose estimation: errors increase
more on unseen participants than on unseen gestures. Second, EMG and vision offer
complementary information. Vision provides the stronger standalone signal in the
current benchmark, whereas EMG captures pose-relevant cues for fine articulation
and improves average accuracy when fused with the evaluated vision baselines.
Third, scaling behavior depends on the data regime. The proposed \model{}
architecture improves over vemg2pose on the original EMG2Pose benchmark, but
larger variants do not consistently improve on \dataset{}. This suggests that
increasing model capacity alone is insufficient under limited subject coverage
and motivates large-scale EMG pretraining, cross-dataset transfer, and
subject-robust adaptation.

\paragraph{Limitations and future work.}
The present study has several limitations. First, \dataset{} includes 53
participants, which is substantial for synchronized multimodal capture but still
limited for modeling the full extent of inter-subject EMG variability. Expanding
the dataset to more diverse participants would enable a more systematic study of
cross-user robustness. Second, the current vision benchmark is limited to single-frame models.
Since \dataset{} provides synchronized multi-view video and RGB-D streams,
it can support future work on temporal, multi-view, and RGB-D-based methods.
The benchmark also does not yet include rich hand-object interaction with instrumented or
densely annotated objects; future extensions should test whether the same
multimodal signals remain helpful when object contact introduces additional
occlusion and pose ambiguity. Third, our fusion baseline is intentionally simple, and its gains are modest for the larger visual backbones, possibly because the growing parameter imbalance between the visual backbone and the compact EMG and fusion branches makes multimodal optimization more difficult. Future work should explore more balanced and expressive fusion architectures, such as cross-attention and uncertainty-aware fusion.

More broadly, \dataset{} provides a testbed for studying when physiological
signals add information beyond visual observations, especially under unreliable
visual evidence such as self-occlusion and fast motion. We hope it will support
future work on subject-robust, multimodal hand-pose estimation under a unified
benchmark.

\section*{Ethics Statement}

All participants provided informed consent under an IRB-approved protocol and were compensated at a standard hourly rate commensurate with local guidelines. Participants were explicitly informed of their right to withdraw from the study at any time and to request deletion of their data; no participant exercised this right. The consent form explained the purpose of data collection, the modalities recorded, and the intended research use.

Direct identifiers are removed from the dataset: facial regions in egocentric and external camera frames are automatically blurred, and participant identifiers are replaced with anonymous numeric IDs. While EMG and motion capture data do not reveal sensitive health information beyond gross motor function, we note that MANO shape parameters and EMG signals could in principle serve as soft biometric identifiers under certain conditions. The dataset will be released for noncommercial research under CC-BY-NC 4.0, accompanied by explicit guidance against re-identification attempts. We commit to hosting the dataset on a stable platform (e.g., Zenodo) with a DOI and versioned releases, and will provide a contact email for data-related inquiries and error reporting.

\section*{Reproducibility Statement}

The anonymized Code and Data Supplement submitted with this manuscript contains training and evaluation code, exact configuration files, normalization statistics, and a runnable numeric reviewer sample. The supplementary material provides the architecture and optimization settings used for the reported \dataset{} experiments. The public release will include dataset documentation, split definitions, preprocessing scripts, benchmark configurations, and an anonymized download link. The reported \dataset{} \model{} variants use 12{,}000-sample (6\,s) target-wrist EMG windows, are trained for 150 epochs with bf16 mixed precision, and use four to six NVIDIA RTX 4090 GPUs depending on model size.

% aaai2027.sty already sets \bibliographystyle{aaai2027} automatically
\bibliography{references}

\ifmainonly
\else
\clearpage
\setcounter{page}{1}
\onecolumn
\section*{Supplementary Material}
\appendix

\begingroup
\small
\setlength{\parskip}{1pt}
\noindent\textbf{Contents}\par
\smallskip
\noindent
\textbf{\ref{app:dataset_details}\quad Dataset details}
  \dotfill \pageref{app:dataset_details}\par
\hspace*{1.5em}\ref{app:gesture_vocabulary}\quad Gesture Vocabulary
  \dotfill \pageref{app:gesture_vocabulary}\par
\hspace*{1.5em}\ref{app:dataset_card}\quad Dataset Card and Release Format
  \dotfill \pageref{app:dataset_card}\par
\hspace*{1.5em}\ref{app:emg_processing}\quad EMG Signal Processing
  \dotfill \pageref{app:emg_processing}\par
\smallskip
\noindent
\textbf{\ref{app:label_validation}\quad Label reconstruction and quality}
  \dotfill \pageref{app:label_validation}\par
\hspace*{1.5em}\ref{app:markers2mano}\quad Markers2MANO Pipeline
  \dotfill \pageref{app:markers2mano}\par
\hspace*{1.5em}\ref{app:wrist_angles}\quad Wrist Articulation Angles
  \dotfill \pageref{app:wrist_angles}\par
\hspace*{1.5em}\ref{app:joint_angle_conversion}\quad MANO-to-UmeTrack Conversion
  \dotfill \pageref{app:joint_angle_conversion}\par
\smallskip
\noindent
\textbf{\ref{app:implementation}\quad Training and implementation details}
  \dotfill \pageref{app:implementation}\par
\hspace*{1.5em}\ref{app:training_details}\quad EMGFormer Training Details
  \dotfill \pageref{app:training_details}\par
\hspace*{1.5em}\ref{app:training_details_vision}\quad Vision and Fusion Training Details
  \dotfill \pageref{app:training_details_vision}\par
\smallskip
\noindent
\textbf{\ref{app:extended_results}\quad Extended experimental results}
  \dotfill \pageref{app:extended_results}\par
\hspace*{1.5em}\ref{app:emg_per_hand}\quad Per-Hand EMG Results
  \dotfill \pageref{app:emg_per_hand}\par
\hspace*{1.5em}\ref{app:per_joint}\quad Per-Joint Error Analysis
  \dotfill \pageref{app:per_joint}\par
\hspace*{1.5em}\ref{app:per_gesture}\quad Per-Gesture Complementarity
  \dotfill \pageref{app:per_gesture}\par
\hspace*{1.5em}\ref{app:occlusion}\quad Self-Occlusion Analysis
  \dotfill \pageref{app:occlusion}\par
\par
\endgroup
\medskip

\section{Dataset details}
\label{app:dataset_details}

\subsection{Gesture Vocabulary}
\label{app:gesture_vocabulary}

Table~\ref{tab:gesture_vocabulary} lists the complete gesture vocabulary of \dataset{} with descriptions.

\begingroup
\scriptsize
\setlength{\tabcolsep}{3pt}
\renewcommand{\arraystretch}{0.98}
\begin{center}
\captionof{table}{Complete gesture vocabulary of \dataset{} (60 classes) with per-gesture descriptions.}
\label{tab:gesture_vocabulary}
\begin{tabular}{@{}p{0.13\textwidth}p{0.31\textwidth}p{0.13\textwidth}p{0.31\textwidth}@{}}
\toprule
\textbf{Gesture} & \textbf{Description} & \textbf{Gesture} & \textbf{Description} \\
\midrule
\multicolumn{2}{l}{\textit{Single-hand (30 gestures)}} &
\multicolumn{2}{l}{\textit{Symmetric bimanual (18 gestures)}} \\
\addlinespace[2pt]
ASL1--ASL9 & American Sign Language digits 1 through 9 &
Clap & Both hands clapping together symmetrically \\
Claw3 & Three-finger claw pose (index + middle + ring flexed) &
CrossHand & Both hands crossed fingers \\
Claw5 & Five-finger claw pose (all fingers flexed) &
CrossStretch & Hands opposite fingers and stretched \\
FreeAction & Free-form unconstrained hand action &
FingerTipTouch & Matching fingertips of both hands touching \\
ILY & ``I Love You'' hand sign (thumb, index, pinky extended) &
FistBump & Two fists bumping together \\
IndexBow & Index finger flexion/extension bowing &
Gaming & Both hands holding a game controller \\
IndexMiddleClaw & Index and middle fingers clawed, others extended &
HandClasp & Hands clasped together \\
JoystickCircle & Circular joystick manipulation on the palm with thumb &
HandRub & Rubbing palms together \\
JoystickSlide & Linear joystick sliding on the palm with thumb &
IndexTapping & Both index fingers tapping in the air \\
MiddleBow & Middle finger flexion/extension bowing &
Kiss & Fingertips touching (thumb and middle), then separating \\
Nine & Hand forming number 9 (four fingers curled, index up) &
PalmStack & One palm stacked on top of the other \\
PalmYaw & Palm facing up/down rotation about the yaw axis &
Prayer & Palms pressed together in prayer position \\
PinchMiddle & Thumb to middle fingertip pinch &
MiddleOppo & Middle and ring fingers of one hand oppose the other hand \\
PinkyBow & Pinky finger flexion/extension bowing &
SymOpen & Both hands opening symmetrically from closed to open \\
Rest & Relaxed hand in neutral resting posture &
SymSwing & Both hands swinging symmetrically side to side \\
RingAndThumb & Ring finger touching thumb tip &
ThumbWrestle & Thumbs wrestling each other \\
RingBow & Ring finger flexion/extension bowing &
Typing & Both hands typing on a virtual keyboard \\
Rock & ``Rock on'' hand sign (index and pinky extended) &
raw & Unconstrained bimanual hand movement \\
Thumb & Thumb up &
\multicolumn{2}{l}{\textit{Asymmetric bimanual (12 gestures)}} \\
\makecell[l]{nocontact\_disperse\\\_palm} & Fingers spread apart with palm open, no contact &
FingerPullLeft & Left hand pulls right-hand fingers \\
nocontact\_free & Free-form finger motion without contact &
FingerPullRight & Right hand pulls left-hand fingers \\
nocontact\_grab & Simulated grasping motion without contact &
PalmRoll & Palms facing, rolling around each other \\
& & PinkyHook & Pinky fingers hooked together \\
& & Squeeze & Both hands squeezing an imaginary object \\
& & Beijing & Two hands crossing to make a ``Bei'' Chinese character \\
& & Checky & Two hands with thumbs forming a ``Checky'' sign \\
& & PairClaw & Both hands clawing with asymmetric finger configurations \\
& & PairOK & Both hands forming OK signs \\
& & Picture & Hands forming a picture frame rectangle \\
& & PinchWring & Two hands pinching and wringing \\
& & ThumbOppo & Thumbs in opposition at different orientations \\
\bottomrule
\end{tabular}
\end{center}
\endgroup

\twocolumn
\subsection{Dataset Card and Release Format}
\label{app:dataset_card}

\paragraph{License.}
The dataset will be released under CC-BY-NC 4.0 for research use, and the baseline code will be distributed under the MIT License. Third-party assets, including the MANO model~\citep{mano} and the training data for \marker{}, remain subject to their original licenses.

\paragraph{File schema.}
Data is organized by episode. Each episode is stored as a parquet file containing synchronized multimodal time series: bilateral EMG (16 channels at 2\,kHz, raw and filtered), per-hand IMU, an egocentric RGB frame from a head-mounted wide-angle GoPro ($1280 \times 720$ at 60\,fps), an external RGB-D frame from a ZED 2i sensor (RGB + depth, $1280 \times 720$ at 30\,fps), MANO pose parameters $\boldsymbol{\theta} \in \mathbb{R}^{48}$ and shape parameters $\boldsymbol{\beta} \in \mathbb{R}^{10}$, wrist flexion/extension and radial/ulnar deviation angles, timestamps, calibration metadata, and split assignments. Video frames are stored as MP4 files and referenced by frame index; each frame includes synchronization diagnostics (\texttt{stale} flag and \texttt{delta\_ms}). Pre-cropped $256\times256$ left- and right-hand patches, generated by reprojecting mocap keypoints onto the egocentric view and extracting per-hand bounding boxes, are provided as per-episode LMDB archives to support vision-based training without requiring users to re-run the projection pipeline. Windows of 12{,}000 EMG samples (6\,s at 2\,kHz) from the target wrist's eight channels are extracted on-the-fly during training.

\paragraph{Access.}
The dataset will be hosted with a download script and checksum verification upon publication.

\subsection{EMG Signal Processing}
\label{app:emg_processing}

All EMG signals are recorded at 2\,kHz. The release includes both raw and filtered EMG for each channel. Filtering is applied per channel in the frequency domain via FFT mask with the following pipeline:
\begin{enumerate}[leftmargin=1.5em,nosep]
    \item Per-channel mean subtraction (DC removal).
    \item Narrow-band notch filters at 50\,Hz (power line frequency) and 100\,Hz (first harmonic).
    \item Soft high-pass transition from 15 to 20\,Hz.
    \item Soft low-pass transition from 850 to 900\,Hz, reaching zero gain at
    900\,Hz.
\end{enumerate}
No amplitude normalization is applied; filtered EMG values are retained in mV. Both raw and filtered streams are preserved in the release. Training uses the filtered EMG by default.

\FloatBarrier
\section{Label reconstruction and quality}
\label{app:label_validation}

\subsection{Markers2MANO Pipeline}
\label{app:markers2mano}

This section details the pipeline used to recover MANO parameters from optical motion capture marker positions during \dataset{} data collection.

\subsubsection{Marker Selection}

We select 21 vertices from the MANO mesh surface as virtual markers, distributed across the hand topology to capture wrist, finger, and palm geometry: indices 191, 88, 253, 708, 729 (wrist/palm), 144, 87, 295, 319 (thumb), 220, 365, 407, 445 (index/middle/ring), and 183, 477, 518, 556, 83, 589, 635, 673 (pinky/palm). These indices are consistent across all MANO shape instances $\boldsymbol{\beta}$, enabling direct correspondence between captured 3D marker positions and the canonical MANO mesh.

\subsubsection{Transformer Architecture}

\begin{center}
\centering
\includegraphics[width=\columnwidth,trim=0 0 0 22,clip]{figures/Transformer_Architecture_aaai.pdf}
\captionof{figure}{Architecture of \marker{} pipeline.}
\label{fig:markers2mano_architecture}
\end{center}

The primary inference model is a Transformer that maps 21 marker positions
$\mathbf{M} \in \mathbb{R}^{21 \times 3}$ to MANO pose
$\boldsymbol{\theta} \in \mathbb{R}^{48}$ and shape
$\boldsymbol{\beta} \in \mathbb{R}^{10}$. The production configuration uses the
500M preset: embedding dimension $d=1600$, $L=16$ pre-norm Transformer layers,
8 attention heads, a $4d=6400$ feed-forward width, dropout 0.1, and
499.87M trainable parameters. The architecture is shown in
Figure~\ref{fig:markers2mano_architecture}. The model consists of:

\begin{itemize}[leftmargin=1.5em,nosep]
    \item \textbf{Input embedding}: Inputs are root-centered using the wrist
    marker. The model accepts all 21 markers, then removes five redundant
    root/MCP markers (indices 1, 5, 9, 13, and 17 in marker order), retaining
    16 tokens. A shared linear layer projects each 3D position to a
    $d$-dimensional token ($\mathbb{R}^3 \rightarrow \mathbb{R}^d$).
    \item \textbf{Transformer encoder}: $L$ pre-norm Transformer layers, each comprising LayerNorm, multi-head self-attention, a residual connection, and a second LayerNorm--feed-forward block with a 4$\times$ expansion ratio ($d \rightarrow 4d \rightarrow d$).
    \item \textbf{Output heads}: After a final LayerNorm, global mean pooling
    feeds an MLP shape head that predicts
    $\Delta\boldsymbol{\beta}\in\mathbb{R}^{10}$. For pose, each retained
    token is concatenated with the global feature and passed through a shared
    MLP that predicts a 6D rotation for one of the 16 MANO joints. The predicted
    pose and shape residuals are added to means computed from valid training
    labels. The heads are zero-delta initialized, and the root rotation is
    fixed to identity. The model does not predict global translation.
    \item \textbf{MANO forward pass}: The predicted $\boldsymbol{\theta}$ and $\boldsymbol{\beta}$ are passed through the MANO layer to obtain the final mesh vertices $\mathbf{V} \in \mathbb{R}^{778 \times 3}$.
\end{itemize}

\subsubsection{Training and Data Augmentation}

The model is trained on hand meshes derived from 11 public datasets using the
augmentation procedure detailed below:
GigaHand~\citep{gigahand}, HOT3D~\citep{hot3d}, ARCTIC~\citep{arctic},
HOI4D~\citep{hoi4d}, HanCo~\citep{hanco}, DexYCB~\citep{dexycb},
H2O~\citep{h2o}, InterHand2.6M~\citep{interhand},
FreiHand~\citep{freihand}, HO3Dv3~\citep{ho3d}, and
HO3Dv2~\citep{ho3d}. Their combined raw footage exceeds 195 million frames.
To balance the training distribution, we assign each dataset a sampling ratio
based on its scale tier: large-scale datasets (100M+ frames, e.g., GigaHand)
at 0.14\%; medium-scale datasets (1M+, e.g., HOT3D, InterHand2.6M) at 1.5\%;
and small-scale datasets (100K+, e.g., FreiHand, HO3D) at 15--20\%. This
yields approximately 0.7M training meshes
(Figure~\ref{fig:markers2mano_data}). This multi-source setup covers a diverse
range of hand shapes, poses, and object interactions.

\begin{figure*}[t]
\centering
\includegraphics[width=0.9\textwidth]{figures/aaai/fig_markers2mano_data_pie.png}
\caption{Composition of the \marker{} training data. Left: original frame counts of 11 public datasets (log scale), spanning 76K to 183M frames. Each dataset is assigned a sampling ratio by scale tier---0.14\% (GigaHand, 100M+), 1.5\% (1M+), or 15--20\% (100K+)---shown as annotations. Right: target distribution of the training mixture after applying the tiered sampling ratios, yielding approximately 0.7M training meshes.}
\label{fig:markers2mano_data}
\end{figure*}

The loss function combines full-mesh supervision with an auxiliary loss on the
16 markers retained by the Transformer:
\begin{equation}
\begin{aligned}
\mathcal{L} ={}&
\frac{1}{V}\sum_{v=1}^{V}
\left\|\mathbf{v}^{\text{pred}}_v-\mathbf{v}^{\text{gt}}_v\right\|_2^2 \\
&+\lambda_{\text{markers}}
\frac{\sum_{n=1}^{16}\alpha_n
\left\|\mathbf{m}^{\text{pred}}_n-\mathbf{m}^{\text{aug}}_n\right\|_2^2}
{\sum_{n=1}^{16}\alpha_n}.
\end{aligned}
\end{equation}
where $\mathbf{v}^{\text{pred}}, \mathbf{v}^{\text{gt}}$ are the predicted and
ground-truth MANO mesh vertices ($V=778$), and
$\mathbf{m}^{\text{pred}},\mathbf{m}^{\text{aug}}$ are the retained predicted
and augmented markers. We set $\lambda_{\text{markers}}=0.5$.
The fingertip weights $\alpha_n$ are 2.0 and all other retained-marker weights
are 1.0.

\paragraph{Optimization.}
We train for 50 epochs with Adam at a learning rate of $5\times10^{-6}$,
batch size 320, bf16 mixed precision, and gradient clipping at 1.0 on five
GPUs. Validation is performed every epoch, and the checkpoint with the lowest
validation loss is retained. The run is initialized from the compatible
warm-start checkpoint specified in the released configuration.

\paragraph{Data augmentation.}
To improve robustness to motion-capture variation, we apply the active
perturbations specified in \texttt{m2m\_config.yaml} to every training batch;
validation inputs are left unperturbed apart from deterministic root centering.

\begin{itemize}[leftmargin=1.5em,nosep]
    \item \textbf{Structural}: 
    (1)~Bone-length perturbation---each bone segment is selected independently
    with probability 0.3 and scaled by up to \(\pm 5\%\) of its length,
    simulating soft-tissue artifacts.
    (2)~Random global scaling---after root-centering, all coordinates are multiplied by a factor in \([0.75, 1.25]\), varying the apparent hand scale.
    \item \textbf{Identity}: 
    (3)~Marker swap---spatially proximate markers (within 15\,mm) are randomly exchanged with probability 0.3 (capped at three swaps per frame).
    (4)~Neighborhood blending---each marker is blended with a random convex combination of its neighboring mesh vertices, retaining a self-weight of 0.5.
    \item \textbf{Noise}: 
    (5)~Gaussian noise---i.i.d.\ \(\mathcal{N}(0,1\,\text{mm}^2)\) noise is added, with the noise suppressed for a marker with probability 0.1.
    (6)~Marker drift---three markers receive a systematic spatial offset of up to 5\,mm in a random direction.
\end{itemize}

Marker dropout and marker-spike augmentation are disabled in the reported
configuration.

To assess the fidelity of markers2mano reconstruction, we report the reconstruction error of the original markers and virtual vertices of the reconstructed meshes. \marker{} shows 3.2\,mm of per-vertex error on the public dataset, and 4.3\,mm of per-marker error on the \dataset{}. We also manually reproject meshes onto egocentric RGB frames sampled randomly to validate the label quality, as shown in Figure~\ref{fig:label_quality}.

\subsection{Wrist Articulation Angle Computation}
\label{app:wrist_angles}

Wrist articulation angles are computed directly from motion-capture markers,
independent of the MANO reconstruction pipeline. Markers attached to each EMG
armband (4 markers on the right wrist, 5 on the left) define a forearm
coordinate frame, from which we compute two wrist degrees of freedom:
flexion/extension $\theta_{\text{fe}}$ and radial/ulnar deviation
$\theta_{\text{ru}}$.

Three non-collinear armband markers define a forearm frame with normalized basis vectors
$\hat{\mathbf{F}} = \mathrm{normalize}(\mathbf{F})$ (forearm direction) and $\hat{\mathbf{L}} = \mathrm{normalize}(\mathbf{L})$ (leftward across dorsal
surface), with normal $\hat{\mathbf{N}} = \mathrm{normalize}(\hat{\mathbf{F}} \times
\hat{\mathbf{L}})$. Let $\hat{\mathbf{H}}$ be the unit vector from wrist to middle
finger MCP, and $\hat{\mathbf{H}}_{\text{proj}} = \hat{\mathbf{H}} - (\hat{\mathbf{H}} \cdot \hat{\mathbf{N}})\hat{\mathbf{N}}$ its projection onto the
$(\hat{\mathbf{F}},\hat{\mathbf{L}})$ plane. The wrist angles are:
\[
\theta_{\text{fe}} = \arcsin(\hat{\mathbf{H}} \cdot \hat{\mathbf{N}}), \qquad
\theta_{\text{ru}} = \operatorname{atan2}\bigl(\hat{\mathbf{H}}_{\text{proj}}
\cdot \hat{\mathbf{L}},\; \hat{\mathbf{H}}_{\text{proj}} \cdot
\hat{\mathbf{F}}\bigr),
\]
with extension and radial deviation defined as positive. For the left hand,
$\mathbf{N} = \mathbf{L} \times \mathbf{F}$ to maintain consistent sign
conventions.

\subsection{Joint Angle Conversion from MANO to UmeTrack / EMG2Pose Format}
\label{app:joint_angle_conversion}

In addition to MANO pose and shape parameters
($\boldsymbol{\theta},\boldsymbol{\beta}$), \dataset{} provides per-frame joint
angles in the 20-DoF scalar format used by EMG2Pose and
UmeTrack~\citep{emg2pose, umetrack}, enabling direct comparison with prior
EMG-to-pose benchmarks that use scalar joint angle representations. The
conversion is formulated as landmark-level inverse kinematics.

\paragraph{Conversion pipeline.}
Given MANO pose and shape parameters for a frame, we first run MANO forward kinematics to obtain 21 joint positions. We align these into the UmeTrack coordinate frame via a fixed transformation (scale, global orientation, and translation) precomputed by jointly optimizing alignment between the two models' rest-pose meshes. We then optimize 20 scalar joint angles $\mathbf{a} \in \mathbb{R}^{20}$ such that the UmeTrack forward kinematics produces landmark positions matching the aligned MANO targets:

\begin{equation}
\mathcal{L}(\mathbf{a}) = \frac{1}{20} \sum_{i=1}^{20} \bigl\|\mathbf{p}_i^{\text{umetrack}}(\mathbf{a}) - \mathbf{p}_i^{\text{mano}}\bigr\|_2^2,
\end{equation}

where $\mathbf{p}_i^{\text{mano}}$ are the aligned MANO joint positions (after applying the fixed scale, rotation, and translation) and $\mathbf{p}_i^{\text{umetrack}}(\mathbf{a})$ are the corresponding UmeTrack landmarks computed via the UmeTrack forward kinematics (scalar joint angles $\to$ axis-angle rotations via Rodrigues $\to$ skinning). The 21 MANO joints are mapped to 20 UmeTrack landmarks via a predefined correspondence; the MANO wrist joint is excluded because it coincides with the UmeTrack coordinate origin and carries no pose information.

\paragraph{Optimization.}
We minimize $\mathcal{L}(\mathbf{a})$ with L-BFGS (strong Wolfe line search,
learning rate 0.1, 100 outer steps with up to 50 inner iterations each).
\begin{samepage}
\noindent Joint range constraints from the UmeTrack hand model are enforced
via sigmoid reparameterization:
\[
\mathbf{a}
= \mathbf{a}_{\min}
+ \left(\mathbf{a}_{\max}-\mathbf{a}_{\min}\right)\odot\sigma(\mathbf{z}),
\]
where $\mathbf{z}$ is the unconstrained optimizable variable and $\odot$
denotes element-wise multiplication. The optimization runs per-episode on GPU
with batch size 256; episodes exceeding GPU memory are processed in smaller
chunks.
\end{samepage}

\paragraph{Output.}
The resulting 22-DoF prediction target combines 20 finger joint angles (indices 0--19) and 2 wrist articulation angles (indices 20--21). The mean per-landmark fitting error of the optimized UmeTrack joints relative to the MANO targets is 2.8\,mm, confirming that the conversion introduces negligible additional error. Table~\ref{tab:joint_angle_semantics} lists the complete semantic ordering per dimension. The conversion script is included in the dataset release.

\begin{center}
\centering
\captionof{table}{Complete 22-DoF joint angle target semantics. Each dimension is reported in degrees. FE = flexion/extension, AA = abduction/adduction. Indices 0--3 cover the thumb, 4--7 the index finger, 8--11 the middle finger, 12--15 the ring finger, and 16--19 the pinky.}
\label{tab:joint_angle_semantics}
\footnotesize
\setlength{\tabcolsep}{2.5pt}
\renewcommand{\arraystretch}{1.02}
\begin{tabular}{clll}
\toprule
Index & Finger & Joint & Motion \\
\midrule
0 & Thumb & CMC & FE \\
1 & Thumb & CMC & AA \\
2 & Thumb & MCP & FE \\
3 & Thumb & IP & FE \\
4 & Index & MCP & AA \\
5 & Index & MCP & FE \\
6 & Index & PIP & FE \\
7 & Index & DIP & FE \\
8 & Middle & MCP & AA \\
9 & Middle & MCP & FE \\
10 & Middle & PIP & FE \\
11 & Middle & DIP & FE \\
12 & Ring & MCP & AA \\
13 & Ring & MCP & FE \\
14 & Ring & PIP & FE \\
15 & Ring & DIP & FE \\
16 & Pinky & MCP & AA \\
17 & Pinky & MCP & FE \\
18 & Pinky & PIP & FE \\
19 & Pinky & DIP & FE \\
20 & -- & Wrist & Flexion/Extension \\
21 & -- & Wrist & Radial/Ulnar deviation \\
\bottomrule
\end{tabular}
\end{center}

We show reconstruction samples in Figure~\ref{fig:label_quality}.

\begin{figure*}[t]
\centering
\includegraphics[width=0.9\textwidth]{figures/recon.png}
\caption{Mesh reprojection and visualization of UmeTrack and MANO mesh reconstruction on \dataset{}. The reconstructed MANO meshes (from \marker{}) are overlaid on egocentric RGB frames, demonstrating accurate hand pose reconstruction and realistic mesh-to-image alignment.}
\label{fig:label_quality}
\end{figure*}

\section{Training and implementation details}
\label{app:implementation}

\subsection{EMGFormer Training Details}
\label{app:training_details}

\paragraph{Optimization.}
For the reported \dataset{} results, all \model{} variants use target-wrist
8-channel EMG windows of 12{,}000 samples (6\,s) and are trained for 150
epochs. We use AdamW~\citep{adamw} ($\beta_1=0.9$, $\beta_2=0.999$, weight
decay $10^{-2}$), an initial learning rate of $5\times10^{-4}$ with cosine
annealing to $5\times10^{-6}$, gradient clipping at 1.0, and bf16 mixed
precision. The objective is L1 joint-angle loss with a fingertip-distance term
of weight 0.01. Per-dataset normalization statistics are computed from the
training data and then fixed for validation and testing. Training uses four to
six NVIDIA RTX 4090 GPUs depending on model size.

\paragraph{Data augmentation.}
We apply the same on-the-fly augmentation family to all three sizes: random
gain, magnitude warping, baseline drift, power-line noise, time/frequency
masking, additive Gaussian noise, channel masking, mixup, and temporal window
jitter. Exact probabilities and ranges are included in the released experiment
configuration.

\paragraph{Architecture and per-size configuration.}
Table~\ref{tab:training_config} summarizes the architecture and training configuration per \model{} variant.

\begin{table}[t]
\centering
\caption{Architecture and training configuration per \model{} variant. ``Params'' includes the TDS encoder, Transformer decoder, and regression head.}
\label{tab:training_config}
\footnotesize
\setlength{\tabcolsep}{2.0pt}
\renewcommand{\arraystretch}{1.05}
\begin{tabular}{lccccc}
\toprule
Variant & $L$ & $d$ & Heads & $d_{\mathrm{FFN}}$ & Params \\
\midrule
\model{}-S  & 3 & 256 & 4  & 512  & 3.5M \\
\model{}-M & 6 & 256 & 8  & 1024 & 6.6M \\
\model{}-L  & 8 & 384 & 12 & 1536 & 16.3M \\
\bottomrule
\end{tabular}

\end{table}

\paragraph{TDS featurizer.}
The TDS-style temporal convolutional frontend~\citep{tds} maps a filtered
2\,kHz EMG window ($T=12{,}000$ samples, 8 target-wrist channels) to a
256-dimensional feature sequence. The architecture consists of two initial
strided Conv1d blocks followed by two TDS stages. Each TDS stage begins with a
strided Conv1d for temporal subsampling, followed by a TDSConvEncoder composed
of alternating depth-separable temporal convolutions (reshaping the channel
axis into a 2D feature map) and channel-wise fully-connected blocks.
Squeeze-and-excitation~\citep{senet} is applied globally after each TDS stage.
The layer-wise configuration is:
\begin{enumerate}[leftmargin=1.5em,nosep]
    \item Conv1d: $8 \to 256$, kernel 11, stride 5
    \item Conv1d: $256 \to 256$, kernel 5, stride 2
    \item TDS Stage 1: in-conv kernel 9, stride 5; 1 block of TDSConv (kernel width 5, channels 8, feature width 32)
    \item TDS Stage 2: in-conv kernel 3, stride 1; 1 block of TDSConv (kernel width 3, channels 8, feature width 32)
\end{enumerate}
The final TDS output is a $(256,230)$ feature map, corresponding to an
effective frame rate of approximately 38\,Hz. The subsequent non-causal
Transformer contextualizes this sequence over the full input window.

\subsection{Vision and Fusion Training Details}
\label{app:training_details_vision}

\paragraph{Vision baseline.}
We provide six fully fine-tuned generic vision backbones: ResNet-\{18, 50, 152\} and ViT-\{Small, Base, Large\}. ResNet backbones are initialized from ImageNet-pretrained weights, while ViT backbones are initialized from DINOv2-pretrained weights. Each generic backbone extracts a feature vector from the window center frame crop ($256\times256$), which is fed through a prediction head to regress 22 joint angles. The prediction head is a 2-layer MLP: $\text{Linear}(d_{\text{backbone}} \to 512) \to \text{ReLU} \to \text{Dropout}(0.1) \to \text{Linear}(512 \to 22)$. All generic vision models are trained with the same L1 (MAE) loss plus fingertip-distance loss (weight 0.01) as the EMG baseline. We additionally fine-tune WiLoR~\citep{wilor} as a hand-specialized baseline using its pretrained MANO-based checkpoint, with the same center-frame supervision target and output dimensionality. Generic backbones use AdamW~\citep{adamw} with cosine annealing and bf16 mixed precision on 6 NVIDIA RTX 4090 GPUs; WiLoR uses AdamW with learning rate $10^{-5}$ and cosine annealing over 30 epochs. Table~\ref{tab:vision_config} lists the per-backbone configuration.

\begin{table}[t]
\centering
\caption{Vision backbone configuration. All models take $256\times256$ center-frame crops as input and are evaluated on the same 22-DoF target. Generic ResNet/ViT models share the same MLP prediction head; WiLoR uses its native hand-pose backbone with a 22-DoF regression head.}
\label{tab:vision_config}
\footnotesize
\setlength{\tabcolsep}{3.0pt}
\renewcommand{\arraystretch}{1.05}
\begin{tabular}{lccccc}
\toprule
Backbone & Params & $d_{\text{backbone}}$ & LR & Batch size & Epochs \\
\midrule
ResNet-18  & 11.5M  & 512  & $10^{-3}$ & 900 & 150 \\
ResNet-50  & 23.5M  & 2048 & $10^{-3}$ & 400 & 150 \\
ResNet-152 & 58.2M  & 2048 & $10^{-3}$ & 180 & 150 \\
ViT-S/14   & 21.9M  & 384  & $10^{-4}$ & 360 & 150 \\
ViT-B/14   & 86.2M  & 768  & $10^{-4}$ & 160 & 150 \\
ViT-L/14   & 303.8M & 1024 & $10^{-4}$ & 50  & 150 \\
WiLoR      & 631.6M & 1280   & $10^{-5}$ & 18  & 30 \\
\bottomrule
\end{tabular}

\end{table}

\paragraph{Fusion baseline.}
The fusion model uses the residual architecture described in the main paper. The vision pathway (backbone, projection layer, and vision head) is initialized from the corresponding vision-only checkpoint, and the EMG pathway (\model{}-S, comprising a TDS featurizer and a Transformer decoder) is initialized from an EMG-only checkpoint. The fusion projection and residual regression head are randomly initialized, with the last layer set to zero so that the initial prediction matches the vision-only baseline ($\Delta\mathbf{y}_{\text{emg}} \approx 0$). The model predicts only the center frame of the window and is trained with AdamW~\citep{adamw}, cosine learning-rate annealing, the same L1 plus fingertip-distance objective, and bf16 mixed precision. We adjust batch size to backbone memory requirements. Exact per-run learning rates, cosine endpoints, epoch budgets, batch sizes, sampling strides, dataset-repeat factors, augmentation settings, and visual-update strategies are provided in the accompanying Hydra configurations.

\FloatBarrier
\section{Extended experimental results}
\label{app:extended_results}

\subsection{Frozen vs. Fine-Tuned Vision in Fusion}
\label{app:fusion_freeze_ablation}

Table~\ref{tab:fusion_freeze_ablation} reports the complete $2\times7$
comparison underlying the model selection in Table~\ref{tab:egoemg_results}.
The paired variants use the same residual-fusion architecture, pretrained
single-modality initializers, and center-frame supervision, while the visual
predictor is either frozen or jointly fine-tuned. Backbone-specific optimization
and sampling settings are recorded in the accompanying experiment configurations.
Selection is performed independently for each backbone using validation MAE.

\begin{table*}[t]
\centering
\caption{Complete frozen-versus-fine-tuned vision ablation for EMG+vision
fusion. Gesture/User/Both report per-subject mean $\pm$ standard deviation in
degrees; Avg is the mean over all splits (lower is better).
Shading marks the variant selected by validation MAE.}
\label{tab:fusion_freeze_ablation}
\footnotesize
\setlength{\tabcolsep}{7.0pt}
\renewcommand{\arraystretch}{1.08}
\begin{tabular}{llcccc}
\toprule
Backbone & Update & Gesture & User & Both & Avg \\
\midrule
ResNet-18 & Frozen & $5.04_{\pm1.25}$ & $6.74_{\pm0.33}$ & $6.64_{\pm0.79}$ & 5.83 \\
\rowcolor{gray!12}
           & Fine-tuned & $4.60_{\pm1.30}$ & $6.40_{\pm1.40}$ & $6.40_{\pm0.80}$ & 5.40 \\
\midrule
\rowcolor{gray!12}
ResNet-50 & Frozen & $4.39_{\pm1.18}$ & $6.19_{\pm0.23}$ & $6.05_{\pm0.60}$ & 5.19 \\
          & Fine-tuned & $4.47_{\pm1.23}$ & $6.20_{\pm0.25}$ & $6.10_{\pm0.59}$ & 5.24 \\
\midrule
\rowcolor{gray!12}
ResNet-152 & Frozen & $4.25_{\pm1.17}$ & $6.04_{\pm0.27}$ & $6.03_{\pm0.59}$ & 5.06 \\
           & Fine-tuned & $4.34_{\pm1.18}$ & $6.11_{\pm0.32}$ & $6.05_{\pm0.64}$ & 5.13 \\
\midrule
ViT-S/14 & Frozen & $5.01_{\pm1.28}$ & $7.33_{\pm0.27}$ & $7.13_{\pm0.32}$ & 6.02 \\
\rowcolor{gray!12}
         & Fine-tuned & $4.50_{\pm1.30}$ & $6.80_{\pm1.20}$ & $6.80_{\pm0.60}$ & 5.54 \\
\midrule
\rowcolor{gray!12}
ViT-B/14 & Frozen & $4.80_{\pm1.26}$ & $7.03_{\pm0.33}$ & $6.84_{\pm0.41}$ & 5.75 \\
         & Fine-tuned & $4.79_{\pm1.27}$ & $7.03_{\pm0.32}$ & $6.87_{\pm0.42}$ & 5.76 \\
\midrule
\rowcolor{gray!12}
ViT-L/14 & Frozen & $4.34_{\pm1.26}$ & $6.69_{\pm0.36}$ & $6.60_{\pm0.56}$ & 5.36 \\
         & Fine-tuned & $4.32_{\pm1.27}$ & $6.69_{\pm0.35}$ & $6.63_{\pm0.57}$ & 5.37 \\
\midrule
\rowcolor{gray!12}
WiLoR & Frozen & $3.96_{\pm1.21}$ & $5.52_{\pm0.42}$ & $5.54_{\pm0.62}$ & 4.68 \\
      & Fine-tuned & $3.95_{\pm1.21}$ & $5.54_{\pm0.39}$ & $5.54_{\pm0.59}$ & 4.69 \\
\bottomrule
\end{tabular}
\end{table*}

\subsection{Per-Hand Breakdown of EMG Baseline Results}
\label{app:emg_per_hand}
Table~\ref{tab:emg_per_hand} reports per-hand MAE for representative methods on \dataset{}, complementing the aggregate results in Table~\ref{tab:egoemg_emg_results}.

\begin{table}[t]
\centering
\caption{Per-hand EMG-to-pose MAE on \dataset{}. MAE is reported in degrees; lower is better. Left and right hands are processed independently with shared model weights, and left-hand EMG is mirrored to the right-hand convention. Rows are grouped by method family.}
\label{tab:emg_per_hand}
\footnotesize
\setlength{\tabcolsep}{3.8pt}
\renewcommand{\arraystretch}{1.08}
\begin{tabular}{llccc}
\toprule
Method & Hand & Gesture & User & Both \\
\midrule
\multicolumn{5}{l}{\textit{\model{} variants}} \\
\midrule
\model{}-S & Left & 11.8 & 15.8 & 16.7 \\
 & Right & 12.6 & 16.5 & 16.7 \\
\midrule
\model{}-M & Left & 11.3 & 15.8 & 16.6 \\
 & Right & 12.0 & 16.4 & 16.9 \\
\midrule
\model{}-L & Left & 11.3 & 15.7 & 16.5 \\
 & Right & 12.4 & 16.6 & 16.8 \\
\midrule
\multicolumn{5}{l}{\textit{Traditional architectures}} \\
\midrule
vemg2pose & Left & 14.7 & 15.7 & 17.2 \\
 & Right & 15.2 & 17.0 & 17.5 \\
\midrule
NeuroPose & Left & 15.9 & 15.4 & 15.7 \\
 & Right & 15.7 & 16.0 & 16.9 \\
\bottomrule
\end{tabular}

\end{table}

Left and right hand MAE agree within 2.0$^\circ$ across all models and splits, confirming that the shared-weight bilateral design and left-to-right mirroring strategy produce symmetric performance. The per-hand consistency supports the validity of processing both hands through a single model.

\subsection{Per-Joint Error Analysis}
\label{app:per_joint}

Table~\ref{tab:per_joint_egoemg} presents the per-finger and per-joint-group mean absolute error (MAE) for \model{} variants on \dataset{}. Across all three model sizes, distal joints and wrist deviation yield the lowest errors (best MAE $10.0^\circ$ for both), whereas the pinky and mid-phalanx groups are the most challenging (best MAE $16.7^\circ$ and $20.7^\circ$, respectively). This pattern is consistent with hand kinematics and sensing difficulty. Distal joints tend to have a more constrained range of motion, making their angles easier to regularize from EMG signals. In contrast, mid-phalanx joints capture a larger portion of finger flexion and therefore exhibit greater motion variability. The higher error on the pinky is also expected: pinky motion is often weaker, more coupled with the ring finger, and less independently represented in wrist-level EMG, making it harder to distinguish from neighboring-finger activation patterns.

\begin{table}[t]
\centering
\caption{Per-joint MAE ($^\circ$) on \dataset{}, left and right hands pooled within each split, then per-sample weighted average across the Gesture, User, and Both splits. Lower is better. Bold indicates the best mean per row.}
\label{tab:per_joint_egoemg}
\footnotesize
\setlength{\tabcolsep}{2.0pt}
\renewcommand{\arraystretch}{1.05}
\begin{tabular}{lccc}
\toprule
Joint group & \model{}-S & \model{}-M & \model{}-L \\
\midrule
Distal & $10.2$ & $10.1$ & $\mathbf{10.0}$ \\
Wrist deviation & $10.8$ & $10.1$ & $\mathbf{10.0}$ \\
Thumb & $12.8$ & $\mathbf{12.3}$ & $\mathbf{12.3}$ \\
Index & $13.4$ & $\mathbf{13.3}$ & $\mathbf{13.3}$ \\
Proximal & $13.7$ & $\mathbf{13.3}$ & $13.4$ \\
Ring & $15.2$ & $14.7$ & $\mathbf{14.6}$ \\
Middle finger & $15.2$ & $14.8$ & $\mathbf{14.8}$ \\
Wrist flexion & $15.7$ & $14.9$ & $\mathbf{14.6}$ \\
Pinky & $17.3$ & $\mathbf{16.7}$ & $17.0$ \\
Mid-phalanx & $21.6$ & $\mathbf{20.7}$ & $\mathbf{20.7}$ \\
\bottomrule
\end{tabular}
\end{table}

For comparison, Table~\ref{tab:per_joint_emg2pose} details the corresponding per-joint MAE on the EMG2Pose dataset. Because EMG2Pose lacks wrist articulation labels, our evaluation here is restricted to finger joints.

\begin{table}[t]
\centering
\caption{Per-joint MAE ($^\circ$) on the EMG2Pose dataset, per-sample weighted average across all three test splits (stage, user, and user+stage). Because EMG2Pose lacks wrist labels, only finger joints are reported. Bold indicates the best mean per row.}
\label{tab:per_joint_emg2pose}
\footnotesize
\setlength{\tabcolsep}{2.5pt}
\renewcommand{\arraystretch}{1.05}
\begin{tabular}{lccc}
\toprule
Joint group & \model{}-S & \model{}-M & \model{}-L \\
\midrule
Proximal & $9.3$ & $8.9$ & $\mathbf{8.5}$ \\
Index & $10.8$ & $10.3$ & $\mathbf{9.8}$ \\
Thumb & $11.1$ & $10.6$ & $\mathbf{9.9}$ \\
Middle finger & $11.1$ & $10.5$ & $\mathbf{10.0}$ \\
Ring & $11.1$ & $10.6$ & $\mathbf{10.1}$ \\
Distal & $13.3$ & $12.8$ & $\mathbf{12.2}$ \\
Pinky & $13.5$ & $12.9$ & $\mathbf{12.3}$ \\
Mid-phalanx & $14.2$ & $13.4$ & $\mathbf{12.6}$ \\
\bottomrule
\end{tabular}
\end{table}

A cross-dataset comparison shows that mid-phalanx and pinky joints are consistently among the most challenging groups across both benchmarks. However, the easiest groups differ between datasets: proximal joints achieve the lowest MAE on EMG2Pose, whereas distal joints and wrist deviation are easiest on \dataset{}. The overall error magnitudes are also higher on \dataset{}---ranging from $10.0^\circ$ to $21.6^\circ$, compared with $8.5^\circ$ to $14.2^\circ$ on EMG2Pose---which likely reflects differences in participant count, gesture diversity, sensor setup, and the inclusion of wrist articulation in \dataset{}'s 22-DoF prediction target.

\subsection{Per-Gesture Modality Complementarity Analysis}
\label{app:per_gesture}

To characterize where EMG provides complementary information beyond vision,
we evaluate all three modalities---vision-only, EMG-only, and
Fusion-RN18+S---on a per-gesture basis across the full test set (all
splits, both hands). Following the Gesture Vocabulary section, we group the
60 gesture classes into
three semantic families: single-hand (30 gestures), symmetric bimanual (18),
and asymmetric bimanual (12). One bimanual gesture (raw, unconstrained) is excluded from this analysis due to its unbounded pose distribution.
Table~\ref{tab:per_gesture_family} reports
sample-weighted per-family averages across all three modalities.

\begin{table}[t]
\centering
\caption{Per-gesture-family modality comparison on \dataset{} (Fusion-RN18+S). Values are sample-weighted averages across all test splits; MAE is in degrees (lower is better). $\Delta = \text{MAE}_\text{vision} - \text{MAE}_\text{fusion}$; positive $\Delta$ indicates fusion improvement. Gesture families follow the Gesture Vocabulary section.}
\label{tab:per_gesture_family}
\footnotesize
\setlength{\tabcolsep}{2.5pt}
\renewcommand{\arraystretch}{1.08}
\begin{tabular}{lrrrrrr}
\toprule
Gesture family & \#G & \#Samp. & Vis. & EMG & Fusion & $\Delta$ \\
\midrule
Single-hand & 30 & 2,406 & 6.0 & 15.0 & 5.5 & $\mathbf{+0.50}$ \\
Symmetric bimanual & 17 & 1,288 & 5.9 & 14.0 & 5.5 & $+0.38$ \\
Asymmetric bimanual & 12 & 584 & 6.8 & 14.6 & 6.4 & $+0.41$ \\
\midrule
Overall & 59 & 4,278 & 6.1 & 14.6 & 5.6 & $+0.46$ \\
\bottomrule
\end{tabular}

\end{table}

Three patterns emerge. First, the fusion gain is broadly distributed: fusion
improves over vision-only on 54 of 59 analyzed gesture classes, suggesting that
the aggregate improvement is not confined to only a few gesture classes.
Second, the largest family-level gain occurs on single-hand
gestures ($\Delta = +0.50^\circ$), followed by asymmetric bimanual
($\Delta = +0.41^\circ$). This suggests that EMG is most useful
when the prediction depends on subtle within-hand articulation or asymmetric
cross-hand coordination that is harder to infer from a single RGB crop.
Symmetric bimanual gestures show the smallest gain ($\Delta = +0.38^\circ$),
as two-hand interactions of this type are often easier to discriminate
visually. Third, even for the few gestures where fusion slightly
underperforms vision-only (Rock, $-0.22^\circ$; Typing, $-0.21^\circ$;
Clap, $-0.12^\circ$; FingerPullLeft, $-0.04^\circ$; CrossHand, $-0.01^\circ$), the degradation is negligible and involves
large-scale hand silhouettes whose visual features are already highly
discriminative. Table~\ref{tab:per_gesture_full} provides the complete
per-gesture breakdown for all 59 analyzed gesture classes, grouped by semantic
family and sorted by descending fusion gain within each family.

\makeatletter
\setlength{\@dblfptop}{0pt}
\makeatother
\begin{table*}[t]
\begingroup
\scriptsize
\setlength{\tabcolsep}{1.8pt}
\renewcommand{\arraystretch}{1.00}
\centering
\captionof{table}{Full per-gesture modality breakdown on \dataset{} using Fusion-RN18+S. MAE is in degrees.
$\Delta = \text{MAE}_\text{vis} - \text{MAE}_\text{fus}$.
Gestures are grouped by family and sorted by descending $\Delta$ within each family. Gesture descriptions are provided in the Gesture Vocabulary section.}
\label{tab:per_gesture_full}
\begin{minipage}[t]{0.48\textwidth}
\centering
\begin{tabular}{rlrrrrr}
\toprule
ID & Name & $N$ & Vis. & EMG & Fus. & $\Delta$ \\
\midrule
\multicolumn{7}{l}{\textit{Single-hand gestures}} \\
6  & ASL7                    & 36  & 7.2 & 16.1 & 6.2 & +0.96 \\
8  & ASL9                    & 38  & 5.7 & 17.5 & 5.0 & +0.71 \\
12 & ILY                     & 40  & 7.3 & 19.2 & 6.6 & +0.70 \\
13 & IndexBow                & 34  & 6.3 & 20.7 & 5.6 & +0.66 \\
20 & PinchMiddle             & 324 & 5.3 & 12.2 & 4.7 & +0.64 \\
26 & Thumb                   & 40  & 7.9 & 24.8 & 7.3 & +0.62 \\
7  & ASL8                    & 42  & 7.0 & 14.8 & 6.4 & +0.61 \\
28 & nocontact\_free          & 292 & 6.8 & 14.2 & 6.2 & +0.61 \\
18 & Nine                    & 320 & 7.0 & 18.1 & 6.4 & +0.59 \\
17 & MiddleBow               & 92  & 4.1 & 11.1 & 3.5 & +0.57 \\
16 & JoystickSlide           & 38  & 5.6 & 19.5 & 5.0 & +0.56 \\
23 & RingAndThumb            & 36  & 7.2 & 16.5 & 6.7 & +0.55 \\
9  & Claw3                   & 38  & 7.7 & 17.3 & 7.2 & +0.54 \\
27 & nocontact\_disperse\_palm & 92 & 4.9 & 12.2 & 4.5 & +0.46 \\
15 & JoystickCircle          & 308 & 4.5 & 12.2 & 4.0 & +0.45 \\
2  & ASL3                    & 30  & 7.3 & 19.6 & 6.8 & +0.44 \\
11 & FreeAction              & 38  & 6.5 & 15.8 & 6.1 & +0.41 \\
5  & ASL6                    & 40  & 6.9 & 18.1 & 6.5 & +0.39 \\
14 & IndexMiddleClaw         & 92  & 5.0 & 12.9 & 4.6 & +0.39 \\
1  & ASL2                    & 34  & 6.9 & 17.7 & 6.5 & +0.38 \\
19 & PalmYaw                 & 82  & 3.2 & 9.8  & 2.8 & +0.36 \\
4  & ASL5                    & 28  & 4.5 & 14.6 & 4.2 & +0.35 \\
24 & RingBow                 & 38  & 7.3 & 16.9 & 7.0 & +0.34 \\
0  & ASL1                    & 44  & 7.9 & 17.6 & 7.7 & +0.26 \\
10 & Claw5                   & 36  & 6.0 & 15.6 & 5.7 & +0.25 \\
22 & Rest                    & 30  & 7.1 & 17.4 & 6.9 & +0.20 \\
21 & PinkyBow                & 40  & 7.5 & 16.6 & 7.4 & +0.18 \\
3  & ASL4                    & 32  & 6.2 & 18.0 & 6.1 & +0.14 \\
29 & nocontact\_grab          & 36  & 7.4 & 18.6 & 7.4 & +0.06 \\
25 & Rock                    & 36  & 9.5 & 17.5 & 9.7 & -0.22 \\
\bottomrule
\end{tabular}
\end{minipage}
\hfill
\begin{minipage}[t]{0.48\textwidth}
\centering
\begin{tabular}{rlrrrrr}
\toprule
ID & Name & $N$ & Vis. & EMG & Fus. & $\Delta$ \\
\midrule
\multicolumn{7}{l}{\textit{Symmetric bimanual gestures}} \\
38 & HandClasp       & 30  & 8.7 & 16.4 & 8.0 & +0.75 \\
52 & MiddleOppo      & 226 & 5.3 & 12.3 & 4.8 & +0.59 \\
39 & HandRub         & 32  & 5.9 & 16.0 & 5.3 & +0.55 \\
36 & FistBump        & 256 & 9.2 & 23.7 & 8.7 & +0.51 \\
44 & Prayer          & 78  & 3.9 & 9.5  & 3.4 & +0.46 \\
42 & PalmStack       & 262 & 4.6 & 10.2 & 4.2 & +0.41 \\
47 & SymSwing        & 28  & 6.6 & 15.7 & 6.2 & +0.36 \\
40 & IndexTapping    & 32  & 5.3 & 13.0 & 5.0 & +0.30 \\
37 & Gaming          & 26  & 7.8 & 15.9 & 7.5 & +0.24 \\
46 & SymOpen         & 30  & 5.0 & 12.6 & 4.8 & +0.22 \\
51 & Kiss            & 78  & 4.0 & 10.4 & 3.7 & +0.21 \\
48 & ThumbWrestle    & 34  & 5.5 & 14.6 & 5.3 & +0.18 \\
35 & FingerTipTouch  & 76  & 3.6 & 7.5  & 3.4 & +0.16 \\
32 & CrossStretch    & 26  & 5.3 & 11.9 & 5.3 & +0.01 \\
31 & CrossHand       & 26  & 5.6 & 12.8 & 5.7 & -0.01 \\
30 & Clap            & 14  & 6.7 & 14.3 & 6.8 & -0.12 \\
49 & Typing          & 34  & 5.3 & 11.9 & 5.5 & -0.21 \\
\addlinespace[3pt]
\multicolumn{7}{l}{\textit{Asymmetric bimanual gestures}} \\
34 & FingerPullRight & 28  & 7.1 & 13.6 & 6.1 & +0.96 \\
58 & PinchWring      & 32  & 7.5 & 17.8 & 6.8 & +0.71 \\
55 & PairClaw        & 36  & 8.1 & 17.4 & 7.4 & +0.66 \\
43 & PinkyHook       & 32  & 7.0 & 16.0 & 6.5 & +0.51 \\
53 & Beijing         & 30  & 10.1 & 19.2 & 9.6 & +0.44 \\
41 & PalmRoll        & 54  & 4.4 & 12.1 & 4.0 & +0.36 \\
59 & ThumbOppo       & 30  & 7.1 & 17.5 & 6.7 & +0.34 \\
45 & Squeeze         & 230 & 6.2 & 12.4 & 5.8 & +0.32 \\
54 & Checky          & 22  & 9.6 & 21.1 & 9.4 & +0.26 \\
56 & PairOK          & 32  & 7.6 & 15.6 & 7.5 & +0.07 \\
57 & Picture         & 32  & 7.1 & 16.7 & 7.1 & +0.03 \\
33 & FingerPullLeft  & 26  & 6.0 & 13.6 & 6.0 & -0.04 \\
\bottomrule
\end{tabular}
\end{minipage}
\endgroup
\end{table*}

\subsection{Fusion Benefit vs.\ Hand Self-Occlusion}
\label{app:occlusion}

To understand when EMG--vision fusion provides the greatest benefit, we
stratify test-set predictions by the degree of hand self-occlusion.

\begin{center}
  \centering
  \includegraphics[width=\columnwidth]{figures/occlusion_fusion_gain.png}
  \captionof{figure}{Fusion benefit versus hand self-occlusion on \dataset{}.
  (a) Vision-only and EMG+vision MAE across 0.05-wide occlusion bins.
  (b) Per-bin fusion gain
  ($\Delta = \mathrm{MAE}_{\mathrm{vision}} - \mathrm{MAE}_{\mathrm{fusion}}$)
  with a LOESS trend. Error bands denote $\pm$1 SEM.}
  \label{fig:occlusion_gain_supp}
\end{center}

\paragraph{Self-occlusion score.}
We quantify per-frame hand self-occlusion via z-buffer rasterization of the
ground-truth MANO mesh in camera space.  Concretely, for each test frame:
\begin{enumerate}[nosep]
  \item The MANO mesh vertices are transformed from world coordinates to camera
        coordinates using the calibrated extrinsic matrix ($\mathbf{v}_{\text{cam}} =
        \mathbf{R}_{C \leftarrow W}\,\mathbf{v}_W + \mathbf{t}_{C \leftarrow W}$).
  \item All mesh triangles are rasterized into a depth buffer (z-buffer) at the
        native video resolution via pinhole projection with the calibrated
        intrinsic matrix~$\mathbf{K}$.  For each pixel covered by a triangle, the
        interpolated depth is stored if it is smaller than the current buffer value.
  \item A vertex is deemed \emph{visible} if, within a $5 \times 5$ pixel
        neighbourhood of its projected location, any z-buffer entry agrees with
        the vertex depth to within $\epsilon = 5$\,mm.
  \item Each triangle's 3-D surface area (half the cross-product norm of two
        edge vectors) is distributed equally to its three vertices, yielding a
        per-vertex area weight~$w_i$.
  \item The \textbf{self-occlusion score} is defined as
        \begin{equation}
          s_{\text{occ}} = 1 - \frac{\sum_{i \in \mathcal{V}_{\text{vis}}} w_i}
                                     {\sum_{i=1}^{V} w_i}\,,
        \end{equation}
        where $\mathcal{V}_{\text{vis}}$ is the set of visible vertices.
        A score of~0 means the hand is fully visible; a score of~1 means it is
        fully occluded from the camera viewpoint.
\end{enumerate}

\paragraph{Results.}
Figure~\ref{fig:occlusion_gain_supp}(a) shows that both vision-only and fusion MAE
vary with occlusion, while the gap between them---the fusion gain---is consistently
positive across the full range, from $\Delta = 0.24\degree$ at the lowest occlusion
bin to $\Delta = 0.52\degree$ at the highest.
The largest fusion benefit occurs at moderate-to-high occlusion
($\Delta \approx 0.8\degree$ near occlusion $\approx 0.72$), where visual evidence
is substantially degraded.
Panel~(b) visualizes the same trend at the per-bin level with a LOESS smooth.
The frame-level association between occlusion and fusion gain is positive but
weak in magnitude (Spearman $\rho = 0.033$, $p < 0.001$; 16{,}650 valid
predictions), so we interpret the bin-wise pattern as suggestive rather than as
evidence of a strong monotonic relationship.

\section*{Supplementary References}
\begingroup
\footnotesize
\setlength{\parindent}{0pt}
\setlength{\parskip}{3pt}

Banerjee, P.; Tuli, S.; Chacko, R.; Kwon, T.; Gupta, A.; and Hebert, M.
2024. HOT3D: A Large-Scale Egocentric Dataset for 3D Hand and Object
Tracking. In \emph{European Conference on Computer Vision (ECCV)}.

Fan, Z.; Taheri, O.; Tzionas, D.; Kocabas, M.; Kaufmann, M.; Black, M. J.;
and Hilliges, O. 2023. ARCTIC: A Dataset for Dexterous Bimanual Hand-Object
Interaction. In \emph{IEEE/CVF Conference on Computer Vision and Pattern
Recognition (CVPR)}.

Hampali, S.; Rad, M.; Oberweger, M.; and Lepetit, V. 2020. HO-3D: A
Benchmark for Hand-Object 3D Pose Estimation. In \emph{IEEE/CVF Conference
on Computer Vision and Pattern Recognition (CVPR)}.

Han, S.; Wu, P.-c.; Zhang, Y.; Liu, B.; Zhang, L.; and Wang, Z. 2022.
UmeTrack: Unified Multi-View End-to-End Hand Tracking for VR. In
\emph{SIGGRAPH Asia Conference Papers}.

Hannun, A.; Case, C.; Casper, J.; Catanzaro, B.; Diamos, G.; Elsen, E.;
Prenger, R.; Satheesh, S.; Sengupta, S.; Coates, A.; and Ng, A. Y. 2019.
Sequence-to-Sequence Speech Recognition with Time-Depth Separable
Convolutions. In \emph{Interspeech}.

Hu, J.; Shen, L.; and Sun, G. 2018. Squeeze-and-Excitation Networks. In
\emph{IEEE Conference on Computer Vision and Pattern Recognition (CVPR)}.

Huang, W.; Guo, H.; Yang, L.; Yao, A.; and Guo, Y. 2024. GigaHand: A
Large-Scale Annotated Dataset for 3D Hand Reconstruction from Egocentric and
Exocentric Views. \emph{arXiv preprint arXiv:2412.04244}.

Kwon, T.; Tekin, B.; St\"uhmer, J.; Bogo, F.; and Pollefeys, M. 2021. H2O:
Two Hands Manipulating Objects for First Person Interaction Recognition. In
\emph{IEEE/CVF International Conference on Computer Vision (ICCV)}.

Liu, Y.; Liu, Y.; Jiang, C.; Lyu, K.; Wan, W.; Shen, H.; Liang, B.; Fu, Z.;
Wang, H.; and Yi, L. 2022. HOI4D: A 4D Egocentric Dataset for Category-Level
Human-Object Interaction. In \emph{IEEE/CVF Conference on Computer Vision
and Pattern Recognition (CVPR)}.

Loshchilov, I.; and Hutter, F. 2019. Decoupled Weight Decay Regularization.
In \emph{International Conference on Learning Representations (ICLR)}.

Moon, G.; Yu, S.-I.; Wen, H.; Shiratori, T.; and Lee, K. M. 2023.
InterHand2.6M: A Dataset and Baseline for 3D Interacting Hand Pose Estimation
from a Single RGB Image. \emph{IEEE Transactions on Pattern Analysis and
Machine Intelligence}, 45(9): 11065--11080.

Potamias, R. A.; Zhang, J.; Deng, J.; and Zafeiriou, S. 2024. WILOR:
End-to-end 3D Hand Localization and Reconstruction in-the-Wild. In
\emph{European Conference on Computer Vision (ECCV)}.

Romero, J.; Tzionas, D.; and Black, M. J. 2017. Embodied Hands: Modeling and
Capturing Hands and Bodies Together. \emph{ACM Transactions on Graphics},
36(6).

Salter, S.; Warren, R.; Schlager, C.; Spurr, A.; Han, W.; Bhasin, A.; Cai,
L.; Walkington, M.; Bolarinwa, M.; Wang, W.; Danielson, C.; Merel, J.;
Pnevmatikakis, E.; and Marshall, L. 2024. emg2pose: A Large and Diverse
Benchmark for Surface Electromyographic Hand Pose Estimation. In
\emph{Advances in Neural Information Processing Systems, Datasets and
Benchmarks Track}.

Yu, B.; Suresh, S.; Mukadam, M.; Maturana, D.; Lih, C.; and Cham, T.-J.
2022. DexYCB: A Benchmark for Enabling Grasp Perception. In
\emph{International Conference on Robotics and Automation (ICRA)}.

Zimmermann, C.; Argus, M.; and Brox, T. 2021. Contrastive Representation
Learning for Hand Shape Estimation. In \emph{DAGM German Conference on
Pattern Recognition}.

Zimmermann, C.; Ceylan, D.; Yang, J.; Russell, B.; Argus, M.; and Brox, T.
2019. FreiHand: A Dataset for Markerless Hand Reconstruction from Single RGB
Images. In \emph{IEEE/CVF International Conference on Computer Vision
(ICCV)}, 10832--10841.

\endgroup


\FloatBarrier
\fi
\end{document}
